\documentclass[11pt]{article}
\usepackage[utf8]{inputenc}
\usepackage[T1]{fontenc}
\usepackage{lmodern}
\usepackage[margin=1in]{geometry}
\usepackage[dvipsnames]{xcolor}
\usepackage{hyperref}
\usepackage{enumitem}
\usepackage{array}
\usepackage{longtable}
\usepackage{booktabs}
\usepackage{tabularx}
\usepackage{titlesec}
\usepackage{parskip}
\usepackage[most]{tcolorbox}
\usepackage{fancyhdr}
\usepackage{setspace}
\hypersetup{
  colorlinks=true,
  linkcolor=BrickRed,
  urlcolor=BrickRed,
  pdftitle={Roman Crime Final Exam Study Guide},
  pdfauthor={OpenAI Codex},
  pdfsubject={Classics 2301 final exam review}
}
\setlength{\parskip}{0.45em}
\setlength{\parindent}{0pt}
\setlength{\headheight}{14pt}
\setlist[itemize]{leftmargin=1.5em, itemsep=0.2em, topsep=0.35em}
\setlist[enumerate]{leftmargin=1.8em, itemsep=0.25em, topsep=0.35em}
\definecolor{RomanRed}{HTML}{8A3B22}
\definecolor{Parchment}{HTML}{F6F0E5}
\definecolor{InkBlue}{HTML}{22313F}
\definecolor{BorderTan}{HTML}{D8C7A4}
\definecolor{SlateBlue}{HTML}{5C6F7B}
\definecolor{OliveNote}{HTML}{5E6D3B}
\definecolor{CompareTeal}{HTML}{315D62}
\titleformat{\section}{\Large\bfseries\color{RomanRed}}{\thesection}{0.6em}{}
\titleformat{\subsection}{\large\bfseries\color{InkBlue}}{}{0em}{}
\newtcolorbox{quotebox}{colback=Parchment,colframe=BorderTan,boxrule=0.6pt,arc=2pt,left=7pt,right=7pt,top=6pt,bottom=6pt}
\newtcolorbox{calloutbox}[1][]{colback=RomanRed!6,colframe=RomanRed!70!black,boxrule=0.8pt,arc=2pt,left=8pt,right=8pt,top=7pt,bottom=7pt,fonttitle=\bfseries,title=#1}
\newtcolorbox{conceptbox}[1][]{colback=SlateBlue!7,colframe=SlateBlue!70!black,boxrule=0.7pt,arc=2pt,left=8pt,right=8pt,top=7pt,bottom=7pt,fonttitle=\bfseries,title=#1}
\newtcolorbox{contextbox}[1][]{colback=OliveNote!8,colframe=OliveNote!70!black,boxrule=0.7pt,arc=2pt,left=8pt,right=8pt,top=7pt,bottom=7pt,fonttitle=\bfseries,title=#1}
\newtcolorbox{comparisonbox}[1][]{colback=CompareTeal!7,colframe=CompareTeal!75!black,boxrule=0.7pt,arc=2pt,left=8pt,right=8pt,top=7pt,bottom=7pt,fonttitle=\bfseries,title=#1}
\newtcolorbox{trapbox}[1][]{colback=BrickRed!6,colframe=BrickRed!70!black,boxrule=0.7pt,arc=2pt,left=8pt,right=8pt,top=7pt,bottom=7pt,fonttitle=\bfseries,title=#1}
\pagestyle{fancy}
\fancyhf{}
\fancyhead[L]{Roman Crime Final}
\fancyhead[R]{\thepage}
\begin{document}
\hypersetup{pageanchor=false}
\begin{titlepage}
\centering
{\Huge\bfseries Roman Crime Final Exam Study Guide\par}
\vspace{1.5em}
{\Large Classics 2301\par}
\vspace{2em}
{\large Crime and Punishment in Ancient Greece and Rome\par}
\vfill
\begin{calloutbox}[How to Use This Packet]
Read the core themes and timeline first, then spend most of your time on Weeks 10--12. Use the earlier weeks to define terms, anchor chronology, and build comparisons. Treat the reading passages as exam evidence, not just illustrations.
\end{calloutbox}
\vfill
{\large Prepared from lecture scripts, slide extracts, study questions, and assigned readings in the local course archive.\par}
\vspace{0.75em}
{\large \today\par}
\end{titlepage}
\hypersetup{pageanchor=true}
\pagenumbering{roman}
\tableofcontents
\newpage
\pagenumbering{arabic}
\begin{calloutbox}[Fast Exam Priorities]
Highest priority: Week 10 Roman law institutions, Week 11 Catiline chronology and the \texttt{SCU}, and Week 12 \texttt{Pro Caelio}, \texttt{vis}, and Cicero's character-based strategy.\newline
Second priority: Week 7 Draco and \texttt{Lysias 1}, Week 8 \texttt{Gorgias}, and Week 6 criminology models.\newline
Foundation priority: Weeks 1--4 for chronology, public/private distinctions, tragedy, and Greek legal vocabulary.
\end{calloutbox}

\section*{Overview}
\addcontentsline{toc}{section}{Overview}


This guide is built from the course materials in this folder: lecture scripts, slide decks, study-question files, the assigned PDFs for \texttt{Lysias 1} and \texttt{In Defense of Marcus Caelius}, and the existing local study-note \texttt{.tex/.pdf} files. The lecture scripts are treated as the main teaching source. The exam is cumulative, but the course materials make clear that Weeks 10-12 deserve the heaviest emphasis.

\section{How to use this guide}

\begin{enumerate}

\item Read the "Core Themes" and "Master Timeline" sections first.

\item Study Weeks 10-12 in detail.

\item Use Weeks 1-4 and 6-8 to make comparisons, define terms, and answer cumulative questions.

\item For reading weeks, memorize both the argument of the text and the professor's interpretation of why the passage matters.

\item Use the study-question answers at the end of each week to self-test.

\end{enumerate}

\section{Guide Legend}

This final guide now follows the same general logic as the midterm notes: it is built for exam review, not just broad understanding. The goal is to make the lecture argument, the core concepts, the likely confusions, and the best comparison language visible in one place.

\begin{calloutbox}[How to Read the Weekly Sections]
\textbf{Five Things} sections flag the identification facts the course keeps returning to: author, title, date, location, and language. \textbf{Key terms} mark the labels, institutions, and legal categories most likely to appear in multiple-choice questions or definition prompts. \textbf{MCQ Trap} warnings mark the places where the lecture says the obvious answer is too simple, too modern, or too literal. \textbf{Passage dossiers} help you turn readings into exam evidence by clarifying who is speaking, what the context is, and why the quotation matters.
\end{calloutbox}

Use the weekly sections in two passes. On the first pass, read the lecture-frame paragraphs and concept explanations so you understand what the week is trying to prove. On the second pass, use the study questions, passage dossiers, glossary, and trap callouts to check whether you can recognize the right answer quickly under exam conditions.

\section{Core Themes You Should See Everywhere}

\subsection{1. Crime is not the same thing in antiquity and modernity}

\begin{itemize}

\item Modern systems sharply separate criminal law and civil law, rely on police and public prosecutors, and use imprisonment as a central punishment.

\item Ancient Athens and Rome distinguished public and private matters, but the line was much narrower and much stranger. Many things modern people call crimes were handled as private wrongs.

\item In classical Athens, homicide was not a public crime in the modern sense. In Rome, many violent acts remained outside a strong state criminal apparatus for a very long time.

\end{itemize}

\subsection{2. Self-help matters}

\begin{itemize}

\item Ancient societies expected households, kin groups, patrons, and private citizens to respond to wrongs.

\item Law often regulated self-help rather than replacing it.

\item Greek homicide law and Roman patron-client practice both show that the state did not monopolize force in the modern way.

\end{itemize}

\subsection{3. Rhetoric shapes criminality}

\begin{itemize}

\item Ancient law-court speeches do not just prove facts. They construct identities.

\item Lysias turns Euphiletos into a modest victim and Eratosthenes into a corrupt seducer.

\item Cicero turns Caelius into a charming young hero and Clodia into a scandalous villain.

\item Sallust frames Catiline in contradictory ways: criminal, aristocrat, revolutionary, enemy, and heroic commander.

\end{itemize}

\subsection{4. Crime is often treated as a matter of character, not just action}

\begin{itemize}

\item Modern positivist criminology looks for causes of criminality in biology, environment, and psychology.

\item Ancient texts often do something similar: they explain crime through ancestry, upbringing, moral weakness, gender, class position, and desire.

\item This matters especially in Weeks 11 and 12.

\end{itemize}

\subsection{5. Insider versus outsider is a recurring structure}

\begin{itemize}

\item Criminals are cast as threats to the community.

\item Punishment often removes them from the community: exile in antiquity, imprisonment in modernity.

\item But the dangerous figure in this course is often an insider: a citizen, aristocrat, householder, wife, husband, politician, or general.

\end{itemize}

\subsection{6. State violence develops over time}

\begin{itemize}

\item In both Greece and Rome, early law coexists with private retaliation.

\item Over time, institutions become more formal.

\item In Rome especially, the state gradually centralizes legitimate violence, and the law on \texttt{vis} becomes a key turning point in that story.

\end{itemize}

\section{How to use ancient/modern comparisons on the exam}

Do not treat the course as a hunt for ``the first version of our system.'' A stronger answer does three things at once: it names the ancient institution on its own terms, explains what work that institution performs in its own society, and only then compares that function to a modern criminal-law equivalent.

\begin{comparisonbox}[Comparison Method]
A weak answer says that a \texttt{quaestio perpetua} is ``like a criminal court'' or that Draco's law is ``like self-defense.'' A better answer says that both ancient and modern systems try to channel violence into recognizable procedures, but they do so with different assumptions about who may initiate action, what counts as public harm, and how much the state controls punishment.
\end{comparisonbox}

\section{Master Timeline}

\subsection{Greek timeline}

\begin{itemize}

\item c. 6500-3000 BCE: Neolithic Greece

\item c. 3000-1100 BCE: Bronze Age

\item c. 1100-750 BCE: Greek Dark Age

\item c. 750-480 BCE: Archaic Period

\item c. 594 BCE: Solon's reforms

\item c. 508 BCE: Cleisthenic reforms at Athens

\item 499-479 BCE: Persian Wars

\item 480-323 BCE: Classical Period

\item 431-404 BCE: Peloponnesian War

\item 469-399 BCE: Socrates

\item 458 BCE: Aeschylus' \texttt{Oresteia}

\item 429-347 BCE: Plato

\item 323-30 BCE: Hellenistic Period

\end{itemize}

\subsection{Roman timeline}

\begin{itemize}

\item 753 BCE: Traditional foundation of Rome

\item 509 BCE: End of monarchy, start of Republic

\item early 5th century BCE: Plebeian secession and Struggle of the Orders

\item 450 BCE: Twelve Tables

\item 367 BCE: Praetorship created; consular hierarchy clarified

\item 287 BCE: \texttt{lex Hortensia} makes plebiscites binding on all Rome

\item 149 BCE: \texttt{lex Calpurnia} creates first \texttt{quaestio perpetua}

\item 82-81 BCE: Sulla's dictatorship and expansion of permanent criminal courts

\item 81 BCE: \texttt{lex Cornelia de sicariis et veneficis}

\item 63 BCE: Catilinarian conspiracy

\item 58 BCE: Clodius' law targets executions without trial; Cicero goes into exile

\item 56 BCE: Trial of Caelius, Cicero delivers \texttt{Pro Caelio}

\item 27 BCE: Augustan settlement

\item 426 CE: Law of Citations

\item 439 CE: \texttt{Codex Theodosianus}

\item 528-534 CE: Justinian's codification, \texttt{Corpus Iuris Civilis}

\end{itemize}

\section{Week 1: Introduction to Crime and Punishment in Ancient Greece and Rome}

\subsection{What Week 1 is really doing}

Week 1 is not just an introduction. It sets the method of the course. The main claim is that modern categories like "crime," "criminal justice," "punishment," "democracy," and even "prison" do not map neatly onto ancient Greece and Rome.

\subsection{Key examples and why they matter}

\begin{itemize}

\item Ned Stark executing a deserter: judgment and punishment are not separated.

\item Salem witch trials: formal procedure can still be unjust.

\item Omar Khadr: one event can generate both criminal and civil consequences.

\item Mike Duffy: public scandal, criminal process, and civil litigation can all diverge.

\end{itemize}

\subsection{Ancient case study: Publius Claudius Pulcher and the sacred chickens}

\begin{itemize}

\item Date: 249 BCE, First Punic War.

\item Claudius ignored bad bird omens before attacking at Drepana.

\item He mocked the sacred chickens, attacked anyway, lost badly, and was prosecuted.

\item The case mixes religion, war, politics, procedure, and punishment.

\item It shows that what counts as a punishable public wrong in Rome includes religious failure, not just modern-style criminal conduct.

\end{itemize}

\subsection{Major concepts}

\begin{itemize}

\item Public case versus private case

\item Criminal law versus civil law

\item Self-help

\item Exile versus imprisonment

\item Parricide and the \texttt{culleus}

\item Foucault and historical discontinuity

\end{itemize}

\subsection{High-yield takeaways}

\begin{calloutbox}[Key Takeaway]

\begin{itemize}

\item Ancient Athens and Rome did distinguish public from private matters, but the categories were narrower than in modern law.

\item Homicide in Athens was not prosecuted the way a modern state prosecutes murder.

\item Ancient prisons were generally temporary holding places, not long-term correctional institutions.

\item Exile is one of the most important ancient punishments.

\end{itemize}

\end{calloutbox}

\subsection{Study questions answered}

\begin{enumerate}

\item What comes to mind when you think about crime and criminals? The course wants you to realize that modern answers are shaped by media, police, courts, and prisons, not by universal truths.

\item Where do we get information about crimes and criminals? Mostly through popular culture, news, legal drama, social media, and highly mediated public stories.

\item How are modern and ancient justice systems similar? Both distinguish serious public wrongs from private disputes, use formal procedures, and try to protect the community.

\item How are they different? Ancient systems rely much more on kinship, self-help, exile, status, and weak state enforcement.

\end{enumerate}

\subsection{Deep study layer}

\begin{calloutbox}[What this lecture is arguing]
Week 1 is not just course introduction. It is a warning against category mistakes. The lecture argues that if you bring modern assumptions about crime, prison, police, prosecution, appeal, and rights directly into Greece and Rome, you will misunderstand almost every later reading. The modern examples are there to show that even modern systems already blur criminal and civil law, emergency power, media spectacle, and political maneuvering.
\end{calloutbox}

\begin{conceptbox}[Concept Bubble: self-help / criminal-civil / exile]
\textbf{Self-help} is the expectation that a wronged person, family, or patron may respond directly rather than waiting for a professional state apparatus. \textbf{Criminal versus civil} is a sorting mechanism that exists in both worlds, but the sorting rules differ sharply. \textbf{Exile versus imprisonment} captures a major punishment difference: ancient societies often protect the community by expelling, fining, disabling, or executing, rather than by long-term incarceration.
\end{conceptbox}

\begin{comparisonbox}[Ancient vs Modern]
The Omar Khadr case blurs citizen, enemy, soldier, and criminal. The Mike Duffy example shows that one factual controversy can produce both criminal and civil proceedings. The Claudius Pulcher story shows that Roman public wrongdoing could combine religion, military disaster, assembly procedure, and retrial-like concerns. The comparison is not that Rome had ``our'' justice system in primitive form, but that both societies use extraordinary procedures differently when the state or community feels endangered.
\end{comparisonbox}

\begin{trapbox}[Exam Trap]
Do not write that ancient Greece and Rome had the same justice system as us, only older. Week 1 is designed to break that assumption.
\end{trapbox}

\section{Week 2: Historical Background and Context}

\subsection{Why this week matters}

Week 2 gives you the chronological map required for the whole course. If you cannot place authors, legal reforms, and political systems in the right order, later weeks become confusing.

\subsection{Greek periods you should know}

\begin{itemize}

\item Neolithic

\item Bronze Age

\item Dark Age

\item Archaic

\item Classical

\item Hellenistic

\end{itemize}

\subsection{Roman periods you should know}

\begin{itemize}

\item Regal Period

\item Republican Period

\item Imperial Period

\item Later Roman Empire

\end{itemize}

\subsection{High-yield anchors}

\begin{calloutbox}[Key Takeaway]

\begin{itemize}

\item Draco and Solon are Archaic, before tragedy and before Plato.

\item Aeschylus belongs to the early Classical period.

\item Socrates and Plato belong to the later Classical period.

\item The Roman law material in Weeks 10-12 belongs mainly to the Republic and early Empire, but codification reaches late antiquity.

\end{itemize}

\end{calloutbox}

\subsection{Study questions answered}

\begin{enumerate}

\item What are the major systems of dating chronology for antiquity? BCE/CE dating plus period labels like Archaic, Classical, Hellenistic, Regal, Republic, and Empire.

\item What are the major periods of Greek history? Neolithic, Bronze Age, Dark Age, Archaic, Classical, Hellenistic.

\item What are the major periods of Roman history? Regal, Republic, Empire, later empire.

\item When did Greek culture emerge in the Mediterranean? In the Bronze Age and especially again after the Dark Age with polis formation and colonization.

\item How did Greek culture develop politically? From small communities and aristocratic structures toward poleis, tyrannies, reforms, and in some places democracy.

\item Which wars shaped Greek and Roman history? Persian Wars, Peloponnesian War, Punic Wars, Roman civil wars.

\item What is meant by "Hellenistic"? The post-Alexander Greek world, after 323 BCE.

\item When did Romans control the Mediterranean? Gradually through the middle and late Republic, especially after the Punic Wars.

\item What connections can you draw between Greek and Roman culture? Roman culture borrows heavily from Greek culture, but Roman legal and political development remains distinct.

\end{enumerate}

\subsection{Deep study layer}

\begin{calloutbox}[What this lecture is arguing]
Week 2 argues that chronology is part of interpretation. You cannot explain a law, a play, a speech, or a conspiracy if you cannot place it inside the right political world. A law from an early city-state, a tragedy from democratic Athens, and a Roman speech from the late Republic do not belong to the same institutional environment.
\end{calloutbox}

\begin{conceptbox}[Concept Bubble: periodization]
\textbf{Periodization} is an interpretive tool, not a fact of nature. It helps you connect ideas of crime and punishment to larger political and social changes. \textbf{Historical context} explains what institutions, anxieties, and audiences a text assumes.
\end{conceptbox}

\begin{comparisonbox}[Ancient vs Modern]
Modern students often imagine legal systems as relatively stable national structures. Antiquity is less stable. Regime change, democratization, aristocratic conflict, conquest, and empire repeatedly alter how wrongdoing is classified and punished.
\end{comparisonbox}

\begin{trapbox}[Exam Trap]
Avoid floating names without periods. Draco, Aeschylus, Socrates, Cicero, and Justinian belong to different institutional worlds.
\end{trapbox}

\section{Week 3: Introduction to Greek Law}

\subsection{What this lecture teaches}

Week 3 gives the vocabulary and institutional structure for Greek law, especially Athenian law. It is one of the densest lectures in the course.

\subsection{Core Greek legal ideas}

\begin{itemize}

\item Earliest Greek dispute resolution often involves arbitration, compensation, and vengeance.

\item Written law becomes important in the Archaic period.

\item Law emerges alongside political reform.

\end{itemize}

\subsection{Key terms}

\begin{itemize}

\item \texttt{nomos}: law, custom, convention

\item \texttt{thesmos}: law laid down

\item \texttt{dike}: often private suit or justice

\item \texttt{graphe}: public prosecution

\item \texttt{Heliaia}: people's court system

\item archons

\item Areopagus

\item boule

\item assembly

\item juries

\item lethal self-help

\end{itemize}

\subsection{Draco and Solon}

\begin{itemize}

\item Draco is associated especially with homicide law.

\item Solon is associated with broader reform and with appeal rights.

\item The important point is not "Draco always imposed death," but that Athenian law formalized when force and homicide were legally permitted or forbidden.

\end{itemize}

\subsection{Private versus public cases}

\begin{itemize}

\item Private cases are brought by the injured party or family.

\item Public cases are brought in the name of the community.

\item This distinction is crucial for understanding why homicide can be legally important without being a public crime in the modern sense.

\end{itemize}

\subsection{Typical Athenian jury trial}

\begin{itemize}

\item No professional judge in the modern sense.

\item Large citizen juries.

\item Timed speeches.

\item No lawyers speaking for you in the modern sense; litigants speak, though they may use speechwriters like Lysias.

\item Heavy dependence on rhetoric.

\end{itemize}

% START WEEK 3 MIDTERM FRAME
\subsection{Lecture frame in plain English}

Week 3 argues that Greek law grows out of custom, arbitration, and feud management rather than a modern state monopoly on prosecution. The procedural topic shows that Athenian institutions distribute legal authority across assemblies, councils, magistrates, and large citizen juries.

This vocabulary is the foundation for the Greek legal and literary weeks that follow. Knowing who does what lets you explain how a case actually moves through the Athenian system. The difference explains why ancient litigants spend so much effort constructing character. Athenian procedure makes persuasive speech part of adjudication itself.

\subsection{Key concepts in paragraph form}

\textbf{nomos.} Law, custom, or accepted social norm. It ties law to shared social practice rather than only to statute. It helps explain why Greek law often shades into morality and custom.

\textbf{Lethal self-help.} Legally bounded private violence used in response to certain serious wrongs. It channels retaliation without eliminating it. It prepares you for Draco's homicide law and the logic behind Lysias 1.

\textbf{graphe versus dike.} A public action versus a more private action tied to the injured party. The distinction tells you whose interest the law recognizes and who gets to bring the case. It is central for explaining why some grave harms are not public crimes in the modern sense.

\begin{trapbox}[MCQ Trap]
Do not assume written law automatically means equality before the law. Do not call the Athenian jury system a simple equivalent of a modern jury trial.
\end{trapbox}
% END WEEK 3 MIDTERM FRAME

\subsection{Study questions answered}

\begin{enumerate}

\item What are the major terms for Greek laws? \texttt{nomos}, \texttt{thesmos}, \texttt{dike}, and \texttt{graphe} are among the most important distinctions.

\item What is meant by lethal self-help? Legally permitted private violence used to defend household order or respond to certain wrongs.

\item What are the main sources for Greek law in the pre-Classical period? Epic, inscriptions, later historians, and surviving statutes cited in speeches.

\item What are the main features of Greek law in the Archaic period? Increasing written law, lawgivers, kin-based retaliation still present, and formal limits on violence.

\item Who are the major reformers at Athens? Draco, Solon, and later Cleisthenes in political organization.

\item What are the main organizations of the Athenian system? Archons, Areopagus, boule, assembly, juries.

\item What is the procedural difference between private and public cases? Who brings the case, whose interest is represented, and sometimes how the case is framed.

\item How does a typical Greek jury trial proceed? Complaint, pretrial stages, speeches, witnesses, jury vote, and sometimes a separate penalty phase.

\end{enumerate}

\subsection{Deep study layer}

\begin{calloutbox}[What this lecture is arguing]
Week 3 argues that Greek law develops out of custom, arbitration, and feud management rather than out of a modern state monopoly on prosecution. The main legal problem is often not simply ``a killing happened,'' but ``how do we stop retaliation from spiraling?''
\end{calloutbox}

\begin{conceptbox}[Concept Bubble: \texttt{nomos} / \texttt{thesmos} / \texttt{dike} / \texttt{graphe}]
\textbf{\texttt{nomos}} ties law to custom and accepted social practice. \textbf{\texttt{thesmos}} marks a move toward fixed, public enactment. \textbf{\texttt{dike}} channels a wrong into a more private action connected to the injured party; \textbf{\texttt{graphe}} treats the wrong as broader and civic. \textbf{Lethal self-help} preserves family honor and enforcement capacity in a world without modern policing, while threatening endless reciprocal violence.
\end{conceptbox}

\begin{comparisonbox}[Ancient vs Modern]
A modern criminal trial usually assumes police investigation, public prosecution, legal professionals, and a judge who controls procedure. Athenian law assumes self-representation, citizen juries, strong rhetorical performance, and a thinner professional legal class. Both seek legitimacy through procedure, but the burden of making the case falls on different people.
\end{comparisonbox}

\begin{trapbox}[Exam Trap]
Do not collapse \texttt{dike} and \texttt{graphe}, and do not say that written law eliminates social inequality. Early written law can publicize rules while still encoding hierarchy.
\end{trapbox}

\section{Week 4: Greek Tragedy of Crime - Aeschylus' Oresteia}

\subsection{The Five Things}

\begin{itemize}

\item Author: Aeschylus

\item Title: \texttt{Oresteia}

\item Date: 458 BCE

\item Location: Athens

\item Language: Greek

\end{itemize}

\subsection{What the trilogy is about}

\textbf{The \texttt{Oresteia} tracks a chain of violence inside the house of Atreus.}

\begin{itemize}

\item Agamemnon sacrifices Iphigenia.

\item Clytemnestra kills Agamemnon.

\item Orestes kills Clytemnestra to avenge Agamemnon.

\item The Furies pursue Orestes.

\item Athena creates a court and transforms vengeance into civic adjudication.

\end{itemize}

\subsection{Why the trilogy matters for the course}

\begin{itemize}

\item It dramatizes a world before stable civic justice.

\item It presents violent revenge as both necessary and destructive.

\item It shows the transition from blood vengeance to courtroom justice.

\item It treats crime as moral, familial, civic, and divine at the same time.

\end{itemize}

\subsection{Major themes}

\begin{itemize}

\item Hospitality

\item Kinship violence

\item Gender inversion and anxiety

\item Divine justice versus human justice

\item Pollution

\item Founding of the court

\end{itemize}

\subsection{High-yield interpretation}

\begin{calloutbox}[Key Takeaway]

\textbf{The trilogy does not simply say "law replaces revenge." It stages a compromise.}

\begin{itemize}

\item older forces of vengeance are not destroyed;

\item they are incorporated into civic order as the Eumenides;

\item justice becomes institutional, but only by accommodating older claims.

\end{itemize}

\end{calloutbox}

% START WEEK 4 MIDTERM FRAME
\subsection{Lecture frame in plain English}

The Oresteia week argues that tragedy is a public arena for thinking about justice, guilt, revenge, and institutional change. The second Oresteia topic argues that Athena's court is a political settlement that transforms vengeance into civic order without abolishing its emotional force.

The trilogy provides the course's clearest mythic model for the move from blood revenge to civic adjudication. It gives you a model for explaining how institutions manage violence instead of magically ending it. The trilogy is useful because it shows what the triumph-of-law story usually hides. Institutional justice often survives by channeling emotion rather than eliminating it.

\subsection{Key concepts in paragraph form}

\textbf{Pollution.} A stain or contagion attached to grave wrongdoing, especially bloodshed. It makes homicide a problem for households, cities, and gods, not just for one victim. It explains why the trilogy treats violence as socially expansive.

\textbf{Furies.} Divine agents of kin-based vengeance. They keep old obligations to blood and family alive. The final settlement matters because Athens must incorporate them, not erase them.

\textbf{Eumenides.} The renamed Furies, now integrated into Athens. They preserve fear, sanction, and continuity between old and new justice. The civic court succeeds because it absorbs old powers into the city.

\begin{trapbox}[MCQ Trap]
Do not say the play simply proves revenge is bad and courts are good; it is about compromise and incorporation. Do not reduce Athena's solution to a simple acquittal.
\end{trapbox}
% END WEEK 4 MIDTERM FRAME

\subsection{Study questions answered}

\begin{enumerate}

\item How does the backstory set up revenge? The house of Atreus is already trapped in inherited violence before the trilogy begins.

\item How does Aeschylus describe life before a justice system? As a world where blood must answer blood and no stable neutral authority resolves conflict.

\item Does Clytemnestra fit expected gender norms? No. She is repeatedly coded as masculine, authoritative, strategic, and politically dangerous.

\item What role does hospitality play? Violations of household trust and guest relationships deepen the moral horror of the crimes.

\item Human versus divine justice? Human judgments are limited and political; divine claims are older, broader, and harder to escape.

\item How does Orestes resolve the conflict? He does not resolve it privately. The case is transferred into a court.

\item How does the trilogy reflect court development? Athena's institution of trial mirrors the move toward civic adjudication.

\item What role do divinities play? They embody competing justice systems.

\item Do punishments fit crimes? Not neatly; each punishment generates further guilt.

\item What do the Eumenides preserve? Order, fear, and the older logic of vengeance, now folded into civic life.

\end{enumerate}

\subsection{Deep study layer}

\begin{calloutbox}[What this lecture is arguing]
Week 4 argues that tragedy is not a break from legal history but another place where Athenians think publicly about violence, guilt, legitimacy, and institutional change. The \texttt{Oresteia} dramatizes the movement from household revenge to civic adjudication, but refuses to pretend that the old logic disappears.
\end{calloutbox}

\begin{conceptbox}[Concept Bubble: pollution / Furies / court]
\textbf{Pollution} marks serious wrongdoing as spiritually and socially contagious. The \textbf{Furies} personify the older order of kin-based vengeance and insist that blood crimes demand response. \textbf{Athena's court} creates a civic forum for deciding blood guilt, but succeeds only by incorporating, not erasing, those older vengeance powers.
\end{conceptbox}

\begin{comparisonbox}[Ancient vs Modern]
Modern states often tell a story in which courts replaced vengeance and made justice rational. The \texttt{Oresteia} is more honest: legal institutions do not abolish anger, fear, or memory. They redirect them. The modern parallel lies in debates over vigilantism, victim rights, restorative justice, and the emotional afterlife of violent crime.
\end{comparisonbox}

\begin{trapbox}[Exam Trap]
Do not say Athena simply proves Orestes innocent in a factual sense. The play is about the creation of authority and legitimacy, not only about evidentiary truth.
\end{trapbox}

\section{Week 5}

There is no lecture-content week in this local course folder. The cumulative guide therefore moves from Week 4 to Week 6.

\section{Week 6: Modern Criminology and Ancient Crime}

\subsection{Why this week matters}

Week 6 gives the theoretical toolkit the professor uses to read the ancient materials. You should know the theory names, but more importantly you should know how they are applied in later weeks.

\subsection{Five approaches to defining crime}

\begin{itemize}

\item Legal: crime is law-breaking.

\item Moral: crime is what is wrong.

\item Social: crime is norm violation.

\item Humanistic: crime is violation of fundamental human rights.

\item Social constructionist: societies construct both crime and criminality.

\end{itemize}

\subsection{Positivism}

\begin{itemize}

\item Biological positivism: Lombroso and the search for physical criminal types.

\item Sociological positivism: environment, neighborhood, and measurable social forces.

\item Psychological positivism: mental and internal causal models.

\end{itemize}

\subsection{Walklate's four explanation models}

\begin{itemize}

\item Rational Choice Theory

\item Social Control Theory

\item Relative Deprivation Theory

\item Hegemonic Masculinity

\end{itemize}

\subsection{Otherness}

\begin{itemize}

\item Criminological other: a person who does not fit the expected image of a criminal.

\item Victimological other: a person who does not fit the expected image of a victim.

\end{itemize}

\subsection{High-yield takeaways}

\begin{calloutbox}[Key Takeaway]

\begin{itemize}

\item Ancient texts can look strikingly modern when they explain criminality through personality, desire, class frustration, or social environment.

\item The course repeatedly asks whether ancient writers focus on the act or on the person.

\end{itemize}

\end{calloutbox}

% START WEEK 6 MIDTERM FRAME
\subsection{Lecture frame in plain English}

Week 6 uses modern criminology as an interpretive toolkit rather than as a template to be imposed mechanically on antiquity. The positivism topic explains why ancient writers can look modern when they explain wrongdoing through character, environment, desire, gender, or social frustration.

These definitions let you explain how ancient texts construct criminality beyond strict legal categories. It gives you a language for analyzing cause, stereotype, and social labeling in later weeks. The continuity lies in construction; the difference lies in style and authority. Both are explanatory projects, but they authorize themselves differently.

\subsection{Key concepts in paragraph form}

\textbf{Social construction of crime.} The view that crime and criminality are made by institutions, norms, and power relations. It shifts attention from timeless essence to labeling and social meaning. It is one of the most useful models for reading Catiline, Clodia, and Euphiletos.

\textbf{Positivism.} A family of approaches looking for observable causes of criminal behavior. It organizes explanation around empirical or quasi-empirical causal models. It helps you compare ancient character narratives with modern causal theories.

\textbf{Criminological other.} A person who does not fit the expected image of offender or victim. It reveals how public stereotypes shape who seems believable or dangerous. This matters for Clodia, Caelius, Catiline, and many ancient defendants.

\begin{trapbox}[MCQ Trap]
Do not treat the legal definition as the only definition that matters in the course. Do not call ancient authors criminologists in the disciplinary sense; the lecture is using analogy, not identity.
\end{trapbox}
% END WEEK 6 MIDTERM FRAME

\subsection{Study questions answered}

\begin{enumerate}

\item How is criminology defined? As a flexible field studying crime, criminality, punishment, victims, and social responses.

\item What are the approaches to defining crime? Legal, moral, social, humanistic, social constructionist.

\item How did criminology develop? Through changing emphasis from law and punishment to scientific, social, and cultural explanation.

\item What is positivism? The search for empirical causes of criminal behavior.

\item What are the three positivist strands? Biological, sociological, psychological.

\item What are Walklate's four models? Rational choice, social control, relative deprivation, hegemonic masculinity.

\item What are the criminological and victimological others? People who do not fit expected public stereotypes of offender or victim.

\end{enumerate}

\subsection{Deep study layer}

\begin{calloutbox}[What this lecture is arguing]
Week 6 argues that modern criminology is not the subject of the course so much as the toolkit for reading the ancient material. The lecture does not claim that Greeks and Romans were criminologists in the modern disciplinary sense. It claims that they often ask the same big questions: what counts as crime, what causes it, who gets labeled criminal, and what punishment is for.
\end{calloutbox}

\begin{conceptbox}[Concept Bubble: social construction / positivism / otherness]
\textbf{Social construction of crime} insists that crime is produced through institutions, norms, and power rather than discovered as a timeless natural fact. \textbf{Positivism} looks for causes of criminality in observable traits, environments, and patterns. \textbf{Criminological and victimological others} explain why some offenders or victims are believed immediately while others are read as unlikely, suspicious, or socially unintelligible.
\end{conceptbox}

\begin{comparisonbox}[Ancient vs Modern]
Modern criminology often claims scientific neutrality, but it still depends on public narratives about who looks criminal. Ancient texts do the same work more openly through stories, metaphors, genealogy, and moral language rather than through expert reports or statistics.
\end{comparisonbox}

\begin{trapbox}[Exam Trap]
Do not treat the legal definition of crime as the only one that matters in this course. Week 6 is there precisely to show why legal definitions are too narrow for the ancient evidence.
\end{trapbox}

\section{Week 7: Athens and Draco's Homicide Law / Lysias 1}

\subsection{The Five Things}

\begin{itemize}

\item Author: Lysias

\item Title: \texttt{On the Murder of Eratosthenes} (\texttt{Lysias 1})

\item Date: probably after 403 BCE, before Lysias' death

\item Location: Athens

\item Language: Greek

\end{itemize}

\subsection{Who is who}

\begin{itemize}

\item Author: Lysias

\item Speaker in court: Euphiletos

\item Defendant: Euphiletos

\item Dead man: Eratosthenes

\end{itemize}

\subsection{Why this case matters}

This week shows how Athenian homicide law works in practice and how rhetoric bends law. Euphiletos must show that the killing was legal, not premeditated murder.

\subsection{The key legal issue}

\textbf{The defense relies on the exception in Draco's homicide law for killing a seducer caught in the act. The most important exam point is this.}

\begin{itemize}

\item the law creates an exception to homicide liability;

\item it does not simply say "adultery is punishable by death."

\end{itemize}

\subsection{Why "seducer" matters}

\textbf{The lecture stresses that seduction is treated as worse than rape because seduction threatens.}

\begin{itemize}

\item household authority,

\item legitimacy of children,

\item inheritance,

\item the integrity of the household.

\end{itemize}

\subsection{Euphiletos' strategy}

\begin{itemize}

\item portray himself as ordinary, modest, gullible, and law-abiding;

\item portray Eratosthenes as a serial corrupter of households;

\item claim the killing was lawful revenge, not murder;

\item imply he acted as the law commanded, even though the law is more limited than that.

\end{itemize}

\subsection{Direct passages from the reading}

\textbf{From \texttt{Lysias 1.pdf}, section 4.}

\begin{quotebox}\itshape

"I believe, gentlemen, that what I have to demonstrate is this: that Eratosthenes seduced my wife and corrupted her ... and that I did not commit this deed for money ... except revenge, as the law allows."

\end{quotebox}

\textbf{Why it matters.}

\begin{itemize}

\item This is the entire defense theory in miniature.

\item Euphiletos reframes the case from homicide to justified household revenge.

\item The phrase "as the law allows" is doing major rhetorical work.

\end{itemize}

\textbf{From \texttt{Lysias 1.pdf}, section 21.}

\begin{quotebox}\itshape

"Make sure, then, no one else finds out about this; otherwise, our agreement will be worth nothing. I expect you to show me them in the act. For I do not need words, but clear evidence..."

\end{quotebox}

\textbf{Why it matters.}

\begin{itemize}

\item Euphiletos tries to prove that he waited for lawful, direct confirmation.

\item It supports his claim that he was not acting out of long-term personal hatred.

\item It also reveals calculated preparation, which is exactly why the speech needs to manage the line between plan and lawful response.

\end{itemize}

\textbf{From \texttt{Lysias 1.pdf}, opening.}

\begin{quotebox}\itshape

"If you employed the same standards for others as you do for your own behaviour, there is not a single one of you who would not be furious..."

\end{quotebox}

\textbf{Why it matters.}

\begin{itemize}

\item Euphiletos invites the jury to identify with him personally.

\item The case is not presented as technical law alone, but as shared masculine household experience.

\end{itemize}

% START WEEK 7 MIDTERM FRAME
\subsection{Lecture frame in plain English}

Lysias 1 turns a homicide prosecution into a contest over whether the dead man, not the defendant, was the real criminal threat. The rhetoric topic shows how Lysias builds Euphiletos as an ordinary, credible man while recasting Eratosthenes as the recognizable type of a dangerous seducer.

It is the clearest Greek example of rhetoric stretching a narrow legal exception into a broad moral license. It shows how legal success can depend on identity construction rather than technical doctrine alone. The difference shows how far the state monopoly on violence has moved. The continuity lies in narrative persuasion; the difference lies in legal framing.

\subsection{Key concepts in paragraph form}

\textbf{Seducer.} The adulterer imagined as a corrupter of household order. It turns sexual misconduct into a broader social and civic danger. The category helps Euphiletos make the killing look legally and morally necessary.

\textbf{Lawful self-help.} Private violence presented as enforcement within legal limits. It lets the defendant appear as an enforcer of household order instead of a murderer. The whole speech depends on the jury accepting this reframing.

\textbf{Logographer.} A professional speechwriter for litigants who still speak in their own voice. It lets rhetoric enter the courtroom while preserving the citizen's formal self-representation. Lysias' artistry works because it feels plain and natural.

\begin{trapbox}[MCQ Trap]
Do not say Draco's law simply authorized killing adulterers in general; the exception is narrower than the speech makes it sound. Do not separate the legal and rhetorical issues; the lecture treats them as inseparable.
\end{trapbox}
% END WEEK 7 MIDTERM FRAME

\subsection{Study questions answered}

\begin{enumerate}

\item What are the Five Things? See above.

\item Who is on trial and who is dead? Euphiletos is on trial for killing Eratosthenes.

\item Who speaks and who wrote the speech? Euphiletos speaks; Lysias wrote it.

\item What matters in Draco's law? The homicide exception for killing a seducer caught in the act.

\item What are the main defense arguments? Eratosthenes seduced the wife, the killing was lawful, there was no prior enmity, and Euphiletos acted with witnesses.

\item How do these fit adultery procedure? The speech exaggerates the law's support for killing and makes lawful response sound mandatory.

\item Why might the defense succeed? Because it aligns with Athenian anxieties about household control and jury identification.

\item How does this connect to Week 6? The speech constructs criminality through character, household role, and social norm rather than bare legal definition.

\item Is the approach similar to modern court practice? In one sense yes, because narrative and character matter; in another no, because the whole legal framework is radically different.

\end{enumerate}

\subsection{Deep study layer}

\begin{calloutbox}[What this lecture is arguing]
Week 7 argues that a homicide trial can turn into a struggle over who gets to count as the real criminal. Euphiletos admits the killing, but Lysias tries to make the jury experience the case not as a murder prosecution but as the lawful destruction of a household predator.
\end{calloutbox}

\begin{conceptbox}[Concept Bubble: logographer / seducer / lawful self-help]
\textbf{The logographer} gives a litigant a professionally crafted voice while preserving the fiction that the citizen speaks for himself. \textbf{The seducer} is not merely sexually immoral; he threatens household authority, legitimacy, inheritance, and masculine control. \textbf{Lawful self-help} lets violence appear as enforcement rather than criminal aggression, provided it stays inside a narrow exception.
\end{conceptbox}

\begin{contextbox}[Passage Context: \texttt{Lysias 1 [1]-[4]}]
Speaker: Euphiletos, performing a defense speech written by Lysias. Audience: Athenian jurors in a homicide court, imagined as male householders. Immediate legal stake: he must turn homicide into justified retaliation against a seducer. Rhetorical job: the opening asks jurors to judge as husbands and fathers before they judge as technicians of law.
\end{contextbox}

\begin{contextbox}[Passage Context: \texttt{Lysias 1 [21]}]
Speaker inside the story: Euphiletos speaking to the enslaved woman who has informed him of the affair. Audience outside the story: jurors listening to his self-portrait as cautious and procedural. Rhetorical job: ``I do not need words, but clear evidence'' makes him sound restrained and evidentiary even while he is privately arranging a lethal confrontation.
\end{contextbox}

\begin{comparisonbox}[Ancient vs Modern]
No modern criminal code gives a husband a general right to kill an adulterer caught in the act. The subtler comparison is that modern trials still depend on stories about credibility, respectability, sexual behavior, and whether a defendant seems like the ``kind of person'' who acted reasonably. Ancient Athens is simply more open about letting those social judgments do legal work.
\end{comparisonbox}

\begin{trapbox}[Exam Trap]
Do not write that Draco's law simply says adultery is punishable by death. The speech wants you to feel that, but the legal exception is narrower and more conditional than the rhetoric suggests.
\end{trapbox}

\section{Week 8: Greek Philosophy of Crime - Plato's Gorgias}

\subsection{The Five Things}

\begin{itemize}

\item Author: Plato

\item Title: \texttt{Gorgias}

\item Date: fourth century BCE, set in the late fifth century

\item Location: Athens

\item Language: Greek

\end{itemize}

\subsection{Main figures}

\begin{itemize}

\item Socrates

\item Gorgias

\item Polus

\item Callicles

\end{itemize}

\subsection{Why rhetoric matters}

Rhetoric is the art of persuasion. In Athens, persuasive speech could shape assemblies and trials. The dialogue asks whether persuasion without truth is power, corruption, or both.

\subsection{Socrates versus Polus}

Polus argues that a successful criminal can be happy if he avoids punishment.

\textbf{Socrates argues.}

\begin{itemize}

\item wrongdoing makes the soul worse;

\item escaping punishment is worse than being punished;

\item punishment is a kind of cure.

\end{itemize}

This is one of the clearest ancient arguments for rehabilitation over simple deterrence.

\subsection{Socrates versus Callicles}

\textbf{Callicles argues.}

\begin{itemize}

\item \texttt{physis} (nature) and \texttt{nomos} (convention/law) are opposed;

\item the strong should dominate the weak;

\item justice is a tool invented by the weak majority;

\item self-indulgence is natural.

\end{itemize}

\textbf{Socrates argues.}

\begin{itemize}

\item justice and self-discipline are real goods;

\item unrestrained desire produces disorder;

\item the predatory outlaw cannot live in stable community.

\end{itemize}

\subsection{High-yield takeaways}

\begin{calloutbox}[Key Takeaway]

\begin{itemize}

\item Gorgias is not a law lecture; it is a lecture on competing definitions of crime and justice.

\item Socrates takes a moral and almost religious view of criminality.

\item Callicles pushes a radical social-constructionist and anti-moral line.

\end{itemize}

\end{calloutbox}

% START WEEK 8 MIDTERM FRAME
\subsection{Lecture frame in plain English}

The rhetoric topic argues that persuasive speech is one of the central mechanisms through which law, politics, and public truth are shaped. The punishment topic stages three rival ways of thinking about wrongdoing: success without punishment, punishment as cure, and power beyond conventional justice.

It gives you the philosophical backdrop for understanding both Lysias and Cicero. It lets you map ancient debate onto modern questions about deterrence, rehabilitation, morality, and social construction. The lecture treats rhetoric as a live modern problem, not an antique curiosity. The ancient argument stays useful because it is about purpose, not just policy.

\subsection{Key concepts in paragraph form}

\textbf{Rhetoric.} The art or technique of persuasive speech. It moves juries, assemblies, and publics toward belief and action. In Athens, rhetorical skill is a route to power and a danger to justice.

\textbf{Sophist.} A teacher of persuasive technique and argument. The sophist professionalizes persuasion as a teachable civic skill. He becomes a symbol of both democratic opportunity and manipulative danger.

\textbf{Punishment as cure.} Socrates' claim that punishment improves the diseased soul of the wrongdoer. It redefines punishment as moral repair rather than only pain or deterrence. This is one of the clearest ancient arguments for a rehabilitative model.

\begin{trapbox}[MCQ Trap]
Do not treat rhetoric as decoration; for Plato and the course it is a technology of power. Do not reduce Socrates, Polus, and Callicles to one-word modern theories and stop there.
\end{trapbox}
% END WEEK 8 MIDTERM FRAME

\subsection{Study questions answered}

\begin{enumerate}

\item What are the Five Things? See above.

\item Who are the main speakers? Socrates, Gorgias, Polus, Callicles.

\item What is rhetoric and why does it matter? Persuasive speech; it matters because it shapes law, politics, and public truth.

\item Who is Gorgias in relation to Socrates? A sophist and teacher of rhetoric, unlike Socrates' philosophical method.

\item Who are Polus and Callicles? Successors in the debate: Polus defends power and success, Callicles defends natural domination and anti-conventional ethics.

\item What criminological theories are visible? Social construction, deterrence, rehabilitation, moral versus legal definitions of crime, hegemonic masculinity.

\item Can the speakers be mapped to modern theories? Yes. Polus resembles deterrence and success-centered thinking; Socrates resembles rehabilitation and moral theory; Callicles resembles radical anti-norm naturalism and social-constructionist suspicion of law.

\end{enumerate}

\subsection{Deep study layer}

\begin{calloutbox}[What this lecture is arguing]
Week 8 argues that rhetoric is not secondary to crime and punishment. It is one of the places where crime and punishment are made intelligible. Plato's \texttt{Gorgias} is therefore part legal theory, part political theory, and part critique of persuasive power.
\end{calloutbox}

\begin{conceptbox}[Concept Bubble: rhetoric / sophist / punishment as cure]
\textbf{Rhetoric} produces belief and moves audiences; in a city where citizens judge cases and vote in assemblies, it is a practical technology of power. \textbf{The sophist} represents professionalized persuasion, useful and profitable but politically dangerous. \textbf{Punishment as cure} lets Socrates imagine justice as moral medicine rather than mere retaliation or deterrence.
\end{conceptbox}

\begin{comparisonbox}[Ancient vs Modern]
The lecture explicitly compares rhetoric to advertising, targeted messaging, social media persuasion, and ``fake news.'' The point is not that sophists were influencers in togas. It is that democratic and mass-media societies remain vulnerable to polished persuasion that can outstrip truth, evidence, or justice.
\end{comparisonbox}

\begin{trapbox}[Exam Trap]
Do not reduce Socrates, Polus, and Callicles to three neat modern theories and stop there. The core issue is the clash between rival visions of power, happiness, and punishment.
\end{trapbox}

\section{Week 10: Introduction to Roman Law}

\subsection{Why this is a very high-yield lecture}

Week 10 is one of the most term-heavy lectures in the course. It covers Roman legal history, magistrates, assemblies, jurists, codification, criminal procedure, and the growth of state authority.

\subsection{Origin story of Roman law}

\begin{itemize}

\item Rome begins, in legend, in 753 BCE.

\item Early Rome is governed by kings.

\item Justice in the regal period is largely custom-based.

\item Priests initially control written law.

\item Family law and religious regulation are early formal areas.

\item Patron-client relations act as a parallel social structure for dispute resolution.

\end{itemize}

\subsection{Core early developments}

\begin{itemize}

\item \texttt{imperium}: power to command

\item \texttt{intercessio}: veto or blocking power

\item Struggle of the Orders

\item tribunes of the plebs

\item Twelve Tables, 450 BCE

\item \texttt{provocatio ad populum}: citizen appeal from magistrate to people

\end{itemize}

\subsection{Magistrates and offices you must know}

\begin{itemize}

\item consuls

\item praetors

\item \texttt{praetor urbanus}

\item \texttt{praetor peregrinus}

\item tribunes of the plebs

\item quaestors

\item aediles

\item censors

\item dictator

\item prorogation, proconsuls, propraetors

\end{itemize}

\subsection{Praetors and legal development}

\begin{itemize}

\item Praetors are central to the justice system.

\item They issue edicts when entering office.

\item Law developed from the edicts is called \texttt{ius honorarium}.

\item This is different from more formal law made by assemblies.

\end{itemize}

\subsection{Roman lawyers and jurists}

\begin{itemize}

\item Early legal specialists are the \texttt{pontifices}.

\item Roman legal work includes:

\item \texttt{ad respondendum} - giving legal advice

\item \texttt{ad agendum} - preparing cases

\item \texttt{ad cavendum} - drafting documents

\item Under Augustus and later emperors, some jurists receive \texttt{ius respondendi ex auctoritate principis}.

\end{itemize}

\subsection{Jurists and schools you should know}

\begin{itemize}

\item Marcus Antistius Labeo

\item Proculian school

\item Masurius Sabinus

\item Sabinian school

\item Gaius

\item Papinian

\item Ulpian

\item Paulus

\item Modestinus

\end{itemize}

\subsection{Codification stages}

\begin{itemize}

\item Law of Citations, 426 CE

\item \texttt{Codex Theodosianus}, 439 CE

\item Justinian's \texttt{Corpus Iuris Civilis}, from 528 CE

\item \texttt{Codex Iustiniani}

\item \texttt{Novellae}

\item \texttt{Institutiones}

\item \texttt{Digest}

\end{itemize}

\subsection{Republican constitution and assemblies}

\begin{itemize}

\item SPQR = Senate and People of Rome.

\item Senate: advisory, but hugely authoritative through \texttt{senatus consulta}.

\item Centuriate assembly:

\item wealth-based;

\item favors the rich;

\item elects consuls, praetors, censors;

\item votes on war;

\item hears appeals in capital cases.

\item Tribal assembly:

\item tribe/residence based;

\item more egalitarian;

\item elects quaestors and curule aediles;

\item votes on \texttt{leges}.

\item Plebeian assembly:

\item plebeians only;

\item elects tribunes and plebeian aediles;

\item votes on \texttt{plebiscita};

\item after \texttt{lex Hortensia}, plebiscites bind all Rome.

\end{itemize}

\subsection{Criminal procedure development}

\begin{itemize}

\item Early criminal action comes from magistrates with \texttt{imperium}.

\item Most cases are still private for a long time.

\item \texttt{quaestio}: ad hoc criminal inquiry/court.

\item \texttt{quaestio perpetua}: permanent court, beginning with \texttt{lex Calpurnia} in 149 BCE.

\item Key permanent courts include:

\item \texttt{repetundae}

\item \texttt{maiestas}

\item \texttt{veneficia et sicarii}

\item \texttt{ambitus}

\item \texttt{vis}

\end{itemize}

\subsection{Why \texttt{lex Cornelia de sicariis et veneficis} matters}

\begin{itemize}

\item 81 BCE

\item Establishes a permanent court for murder and poisoning.

\item Also criminalizes carrying weapons with intent to kill or steal.

\item This is one place where you can argue that murder becomes criminal in Rome in a stronger sense than earlier.

\end{itemize}

\subsection{Later empire and class}

\begin{itemize}

\item \texttt{cognitio extraordinaria} eventually takes over much judicial work.

\item Urban prefects and other prefects extend enforcement.

\item \texttt{honestiores} receive better treatment than \texttt{humiliores}.

\item Augustus adds moral legislation, including adultery courts.

\end{itemize}

\subsection{Delicts}

\begin{itemize}

\item Not quite crimes, not quite torts.

\item Core delicts:

\item \texttt{furtum}

\item \texttt{rapina}

\item \texttt{damnum iniuria datum}

\item \texttt{iniuria}

\end{itemize}

% START WEEK 10 MIDTERM FRAME
\subsection{Lecture frame in plain English}

The origins topic argues that Roman law emerges through legend, custom, priestly control, patronage, and political struggle rather than appearing as a ready-made code. The magistrates topic argues that Roman law works through offices, assemblies, and checks on power rather than through a simple separation of branches. The jurists topic argues that Roman law becomes increasingly shaped by expert interpretation, teaching, and codification. The criminal-procedure topic argues that Rome slowly develops more regular public courts for recurring political and social dangers.

It shows that Roman legal history is inseparable from Roman constitutional history. You need these institutions to understand how legislation, appeal, and legal administration operate in the Republic. It explains why Roman law has such a long later legacy even after the political world that produced it has changed. It is the bridge between general Roman legal history and the specific Roman crime material in Weeks 11 and 12. Roman law is foundational, but not because it already looks modern in every respect. Roman law must be explained as a political process, not just a doctrinal one. Roman law's legacy lies heavily in reasoning, classification, and textual authority. The Roman development matters because it is partial, contested, and status-laden.

\subsection{Key concepts in paragraph form}

\textbf{Patron-client system.} A hierarchy in which patrons protect and assist dependents or clients. It channels disputes, protection, and influence through social hierarchy. It modifies self-help and weak state enforcement instead of eliminating them.

\textbf{Twelve Tables.} The early public posting of Roman law in the fifth century BCE. They reduce priestly secrecy and make law more publicly knowable. They are a landmark in access and publicity, not the final form of Roman law.

\textbf{imperium.} The power to command held by certain Roman magistrates. It authorizes coercion, command, and political action. It explains why legal administration is bound to office-holding.

\textbf{intercessio.} The power to block or forbid another magistrate's action. It checks concentration of power inside a collegial magistracy. It is one of Rome's key anti-monarchical mechanisms.

\textbf{provocatio.} Citizen appeal from coercive magistrate action to the people. It limits state coercion against citizens. It becomes crucial background when emergency execution without trial is later debated.

\textbf{ius honorarium.} Law and remedies developed through magistrates, especially praetors, in office. It adapts rigid older law to new disputes and social realities. It shows Roman law's flexibility and practical problem-solving.

\textbf{Codification.} The collection and ordering of legal materials into more authoritative bodies of text. It stabilizes and transmits legal knowledge across time. This is why the later Roman legal tradition matters far beyond the Republic.

\textbf{quaestio perpetua.} A permanent criminal court for a recurring category of wrong. It regularizes prosecution and classification for repeated public harms. It marks Rome's movement toward a more recognizable criminal-law structure.

\textbf{vis.} Unauthorized force or violence treated as a public danger. It shifts violence from a private quarrel frame toward a state-protection frame. It becomes crucial for the Catiline and Caelius material.

\begin{trapbox}[MCQ Trap]
Do not write as if the Twelve Tables are a complete rational code or as if early Rome already had a full public criminal system. Do not call the praetor a straightforward modern judge or the Senate a parliament. Do not reduce codification to memorizing dates; the point is the growth of legal authority through texts. Do not assume all serious Roman violence was always handled as a modern public crime.
\end{trapbox}
% END WEEK 10 MIDTERM FRAME

\subsection{Study questions answered}

\begin{enumerate}

\item \textbf{What are the main features of the origin story of Roman law?} Roman law begins as a layered story rather than as a finished code. The lecture starts with kings, custom, priestly control, household authority, and patronage, then shows how plebeian struggle, publicity, magistracies, and writing slowly make law more visible and more civic.

\item \textbf{Why are the Twelve Tables important?} The Twelve Tables matter because they reduce priestly secrecy and make law publicly knowable. The lecture treats them as a landmark in access and publicity, not as a complete rational code that solved Roman legal development once and for all.

\item \textbf{Which officials are mainly responsible for legal affairs?} In Republican Rome the praetors matter most for day-to-day legal administration, especially the praetor urbanus and praetor peregrinus. That answer should still be placed inside a wider constitutional setting that includes magistrates, assemblies, jurists, and later imperial authority.

\item \textbf{What are the main sources of law?} The lecture wants you to think historically here. Early Rome emphasizes custom, priestly expertise, and the Twelve Tables; the Republic adds leges, plebiscita, praetorian edicts, and senatorial influence; the Empire gives growing weight to jurists and imperial constitutions.

\item \textbf{How did legal education develop?} Roman legal expertise begins in the hands of the pontifices, then becomes a prestigious juristic culture in which elite experts teach, interpret, and classify law. By late antiquity, legal knowledge is increasingly textual, school-based, and tied to state-backed authority.

\item \textbf{What are the main stages in codification?} The course emphasizes three big codifying moments: the Law of Citations in 426 CE, the Codex Theodosianus in 439 CE, and Justinian's Corpus Iuris Civilis from 528 onward. The key point is not just memorizing dates but seeing the long shift from scattered expertise toward curated bodies of authoritative text.

\item \textbf{How did legislation work?} Roman legislation is political before it is purely doctrinal. Magistrates propose, assemblies vote, the Senate advises with enormous practical authority, and later emperors dominate the process more directly.

\item \textbf{How did criminal law and trials develop?} The lecture presents Roman criminal law as gradual and uneven: magistrate coercion comes first, ad hoc quaestiones follow, and permanent courts eventually appear for recurring public wrongs such as extortion, treason, bribery, and violence. That development matters because it marks Rome's slow movement toward stronger public control over harms that were earlier treated more privately or more politically.

\item \textbf{What role did the emperor play?} The emperor becomes increasingly central as law develops: he appoints officials, shapes appeals, strengthens juristic authority, and ultimately stands behind late codification. A strong answer explains that this is part of the wider centralization of interpretive and coercive power.

\end{enumerate}

\subsection{Deep study layer}

\begin{calloutbox}[What this lecture is arguing]
Week 10 argues that Roman law is really a history of institutions, political struggle, and state-building. It is not a single finished code. It is a long process in which custom, priestly knowledge, aristocratic office, plebeian agitation, praetorian innovation, juristic expertise, and imperial authority gradually produce a more formal legal order.
\end{calloutbox}

\begin{conceptbox}[Concept Bubble: \texttt{imperium} / \texttt{intercessio} / \texttt{provocatio} / \texttt{quaestio perpetua} / \texttt{vis}]
\textbf{\texttt{imperium}} is command power. \textbf{\texttt{intercessio}} checks that power by allowing a colleague or tribune to block action. \textbf{\texttt{provocatio}} limits coercion by letting a citizen appeal. \textbf{\texttt{quaestio perpetua}} regularizes recurring public harms as recognizable courts. \textbf{\texttt{vis}} helps reclassify unauthorized force as a threat to public order and the commonwealth.
\end{conceptbox}

\begin{comparisonbox}[Ancient vs Modern]
Roman law is often called an ancestor of modern law, but that is truest for private law and legal reasoning, not for a modern police-centered criminal state. Rome shares the desire to classify harms and formalize procedure, but it differs in its dependence on status, patronage, office-holding, and the weak separation between law and politics.
\end{comparisonbox}

\begin{trapbox}[Exam Trap]
Do not call the praetor a modern judge without qualification, do not call the Senate a legislature in the parliamentary sense, and do not treat the Twelve Tables as a complete rational code. The lecture presents them as a landmark in publicity and access.
\end{trapbox}

\section{Week 11: Criminals and Enemy Combatants - Sallust's Catiline}

\subsection{The Five Things}

\begin{itemize}

\item Author: Sallust

\item Title: \texttt{Catiline's War} (\texttt{Bellum Catilinae})

\item Date: c. 44-40 BCE

\item Location: Rome

\item Language: Latin

\end{itemize}

\subsection{Why this week matters}

This is the course's big Roman case study in how rhetoric, political crisis, and criminality overlap. The central question is whether Catiline is a criminal conspirator or an enemy combatant in a civil war.

\subsection{Background to the conspiracy}

\begin{itemize}

\item Catiline is a noble but indebted aristocrat.

\item Roman politics is split between \texttt{optimates} and \texttt{populares}.

\item Debt crisis matters.

\item Sulla's civil wars and veteran settlements leave instability in Italy, especially Etruria.

\item Gaius Manlius becomes crucial in Etruria.

\end{itemize}

\subsection{Core chronology of 63-62 BCE}

\begin{itemize}

\item 64 BCE: Catiline loses bid for the consulship of 63.

\item 63 BCE: Catiline loses again for the consulship of 62.

\item July 63 BCE: failed election deepens his crisis.

\item October 20: anonymous letters reveal conspiracy.

\item October 21: Senate passes the \texttt{senatus consultum ultimum}.

\item October 27: Manlius raises revolt in Etruria.

\item early November: Catiline charged with \texttt{vis}.

\item November 7: attempted assassination of Cicero fails.

\item November 8: Cicero delivers the First Catilinarian.

\item mid-November: Catiline declared \texttt{hostis}.

\item December 3: Allobroges produce evidence; conspirators arrested.

\item December 5: Senate advises execution of captured conspirators without trial.

\item January 62 BCE: Catiline dies in battle.

\end{itemize}

\subsection{Why the legal issue is so important}

The captured conspirators were Roman citizens. They should still have had \texttt{provocatio}. Cicero executes them without trial under cover of the \texttt{SCU}, and that choice later contributes to his exile.

\subsection{Criminological reading of Sallust}

\textbf{Sallust does not present Catiline simply as a villain. He mixes.}

\begin{itemize}

\item psychological explanation,

\item moral corruption,

\item class resentment,

\item elite decadence,

\item relative deprivation,

\item rational calculation,

\item charisma and heroism.

\end{itemize}

\subsection{Direct passages from the reading as quoted in the lecture script}

\textbf{From the Week 11 lecture, quoting Sallust on the urban plebs.}

\begin{quotebox}\itshape

"the entire plebs, in its enthusiasm for revolution, approved completely of Catiline's projects"

\end{quotebox}

\textbf{Why it matters.}

\begin{itemize}

\item Catiline is not a lone deviant.

\item The crisis has a social base.

\item This supports a relative-deprivation reading of the conspiracy.

\end{itemize}

\textbf{From the Week 11 lecture, Sallust on Crassus.}

\begin{quotebox}\itshape

"everyone cried that the informant was false ... while very many of them were beholden to Crassus through personal business"

\end{quotebox}

\textbf{Why it matters.}

\begin{itemize}

\item Status shapes legal truth.

\item Roman political justice is shown as selective and class-bound.

\item Power can move a man out of the category of suspect altogether.

\end{itemize}

\textbf{From the Week 11 lecture, Caesar's proposal.}

\begin{quotebox}\itshape

"their money should be confiscated and they themselves held in chains in the municipalities..."

\end{quotebox}

\textbf{Why it matters.}

\begin{itemize}

\item Caesar argues against execution without trial.

\item He implicitly proposes long-term detention, which is strikingly unclassical.

\item He treats the captives more like criminals than battlefield enemies.

\end{itemize}

\textbf{From the Week 11 lecture, Cato's deterrence argument.}

\begin{quotebox}\itshape

"when you decide about P. Lentulus and the others, be assured that at the same time you are issuing a decree about Catiline's army"

\end{quotebox}

\textbf{Why it matters.}

\begin{itemize}

\item Cato shifts the frame from prosecution to war.

\item He justifies extreme punishment as preemption and deterrence.

\item This is the decisive move from "criminal conspiracy" to "enemy threat."

\end{itemize}

\subsection{Catiline's contradictory image}

\textbf{Sallust makes Catiline.}

\begin{itemize}

\item depraved but charismatic,

\item criminal but aristocratic,

\item socially dangerous but politically comprehensible,

\item enemy of Rome but heroic in battle.

\end{itemize}

At the end he dies bravely, which keeps the category unstable.

% START WEEK 11 MIDTERM FRAME
\subsection{Lecture frame in plain English}

The chronology topic shows that the Catilinarian crisis unfolds as a rapidly escalating sequence of elections, plots, emergency decrees, arrests, executions, and battle. The debate topic argues that Caesar and Cato are really debating whether the captured conspirators remain punishable citizens or have become war enemies.

Knowing the order of events is essential because the legal arguments depend on when the conspirators are still citizens and when they are treated as enemies. This is the course's clearest Roman case study in the criminal-versus-enemy problem. The chronology matters because emergency narratives build step by step. The continuity lies in the pressure emergency politics places on ordinary procedure.

\subsection{Key concepts in paragraph form}

\textbf{SCU.} The senatus consultum ultimum, the Senate's ultimate emergency decree. It authorizes extreme action in the name of saving the state while leaving legality contested. It is the background for Cicero's later vulnerability and exile.

\textbf{hostis.} Enemy, in a politically and legally charged sense. It pushes a dangerous insider outside ordinary civic protection. The category makes exceptional punishment feel easier to justify.

\textbf{Relative deprivation.} A model explaining crime or rebellion through blocked expectations and frustrated status. It ties support for Catiline to structural grievance rather than isolated wickedness. It helps explain why Sallust presents social causes alongside moral corruption.

\begin{trapbox}[MCQ Trap]
Do not collapse the whole crisis into one moment; the legal and rhetorical framing changes as the timeline advances. Do not turn Caesar into a pure civil-liberties hero or Cato into mere bloodthirst; both are arguing from public-safety logic, but with different frames.
\end{trapbox}
% END WEEK 11 MIDTERM FRAME

\subsection{Study questions answered}

\begin{enumerate}

\item \textbf{What are the Five Things?} For the exam, the basic frame is Sallust, Catiline's War, written around 44-40 BCE, set in Rome, and written in Latin. The point of memorizing the Five Things is to anchor the reading in late Republican political crisis rather than treating it as free-floating moral literature.

\item \textbf{What social issues lead Catiline toward revolution?} The lecture highlights debt, blocked aristocratic ambition, veteran unrest, and the aftereffects of civil conflict. Sallust therefore presents conspiracy not only as personal wickedness but also as something that feeds on structural grievance and social instability.

\item \textbf{What is happening in Etruria?} Etruria is the crucial armed base of the crisis, especially under Manlius. On the exam, treat it as the place where political conspiracy starts to become something closer to open revolt or civil war.

\item \textbf{Who is Catiline?} Catiline is an indebted noble, a failed candidate, a charismatic political actor, and an unstable criminal category all at once. Sallust keeps him morally corrupt, socially intelligible, and even heroic in battle, which is why the lecture insists that he cannot be reduced to a one-word label.

\item \textbf{Who is Cicero?} Cicero is consul, orator, and the self-fashioned savior of the Roman state. He is also the figure whose later vulnerability shows how dangerous emergency punishment without trial could be, even when it was done in the name of public safety.

\item \textbf{How does Sallust explain Catiline's criminality?} Sallust mixes moral corruption, ambition, desire, elite decadence, social crisis, and political grievance. That mixture is the point: Catiline is explained at once through character and through the diseased condition of the Republic.

\item \textbf{How does Catiline explain his own turn to crime?} In the lecture's terms, Catiline sounds like a figure of grievance, exclusion, and frustrated entitlement. That makes him useful for comparison with relative deprivation and rational-choice style explanations, even though Sallust never stops framing him as dangerous.

\item \textbf{How do Caesar and Cato relate to crime and punishment?} Caesar resists execution without trial and warns about precedent, but he still imagines harsh detention and control. Cato shifts the frame toward deterrence, war, and political signaling, which is why the lecture treats their disagreement as a debate over whether the captives are still citizens or already enemies.

\item \textbf{What is the outcome and lasting issue?} The conspiracy is crushed and Catiline dies in battle, but the legal problem remains unresolved. The lasting issue is when a state may stop treating insiders as citizens subject to procedure and begin treating them as enemies who can be neutralized exceptionally.

\end{enumerate}

\subsection{Deep study layer}

\begin{calloutbox}[What this lecture is arguing]
Week 11 argues that the Catilinarian crisis is not only a story about one conspiracy. It is a case study in what happens when a state cannot decide whether to treat dangerous insiders as citizens, criminals, rebels, or enemies. The whole lecture is framed through modern terrorism discourse in order to show how emergency politics pressures law.
\end{calloutbox}

\begin{conceptbox}[Concept Bubble: \texttt{SCU} / \texttt{hostis} / relative deprivation]
\textbf{The \texttt{SCU}} authorizes extreme emergency action in the name of saving the state while leaving legality contestable. \textbf{\texttt{hostis}} pushes a citizen toward the category of enemy and makes extraordinary violence feel necessary. \textbf{Relative deprivation} explains why Catiline can attract support from debtors, displaced people, and politically frustrated elites rather than appearing as a lone deviant.
\end{conceptbox}

\begin{contextbox}[Passage Context: Sallust on the plebs / Crassus / Caesar / Cato]
Sallust as narrator uses the plebs and Crassus passages to widen the crisis beyond one man and to show how status shapes credibility. Caesar, speaking to the Senate on 5 December 63 BCE, warns against irreversible precedent and tries to keep the captives inside a punitive rather than battlefield frame. Cato, in the same debate, turns sentencing into a signal to Catiline's army and makes deterrence and preemption do the work of justification.
\end{contextbox}

\begin{comparisonbox}[Ancient vs Modern]
The lecture invites comparison to homeland-security logic, preventive detention, covert surveillance, and the claim that terrorism is too dangerous to wait for completed action. The strongest comparison is not that Catiline was exactly a modern terrorist, but that both ancient and modern states feel pressure to move from proof and trial toward suspicion, preemption, and exceptional punishment when leaders can plausibly say the whole community is at risk.
\end{comparisonbox}

\begin{trapbox}[Exam Trap]
Do not turn Caesar into a modern civil-liberties saint, and do not say that the \texttt{SCU} simply made executions legal. The controversy exists precisely because legality remained unstable.
\end{trapbox}

\section{Week 12: Character and Ciceronian Crime - Pro Caelio}

\subsection{The Five Things}

\begin{itemize}

\item Author: Marcus Tullius Cicero

\item Title: \texttt{In Defense of Marcus Caelius} (\texttt{Pro Caelio})

\item Date: 56 BCE

\item Location: Rome

\item Language: Latin

\end{itemize}

\subsection{Background to the case}

\begin{itemize}

\item Caelius is charged with \texttt{vis}, political violence.

\item Prosecutors: Atratinus, Publius Clodius, and Herennius Balbus.

\item Defense: Caelius, Crassus, and Cicero.

\item The wider background involves Egyptian politics, Ptolemy XII, the Alexandrian embassy, and the death of Dio.

\item Clodia, Caelius' former lover, becomes central to Cicero's version of events.

\end{itemize}

\subsection{Why the charge matters}

\texttt{Vis} is supposed to cover dangerous violence against the state, not private scandal. Cicero's whole strategy is to make the jury think the charge is absurdly overblown.

\subsection{Cicero's main strategy}

\textbf{He does not focus primarily on reconstructing facts. He.}

\begin{itemize}

\item turns the trial into comedy;

\item turns Caelius into a charming young hero;

\item turns Clodia into the scandalous manipulator behind the prosecution;

\item makes the case about character instead of action.

\end{itemize}

\subsection{Direct passages from the reading}

\textbf{From \texttt{Cicero In Defense of Marcus Caelius.pdf}, opening section.}

\begin{quotebox}\itshape

"If somebody were here who knew nothing about our laws, our courts, and our customs, he would certainly wonder what crime could be so vicious ... this one court should be in session."

\end{quotebox}

\textbf{Why it matters.}

\begin{itemize}

\item Cicero opens by exaggerating the seriousness of the legal frame.

\item He then undercuts it by suggesting the charge of \texttt{vis} is wildly disproportionate to the real case.

\item He prepares the jury to see the prosecution as theatrical overreach.

\end{itemize}

\textbf{From \texttt{Cicero In Defense of Marcus Caelius.pdf}, on the charges.}

\begin{quotebox}\itshape

"There are two charges. One involves gold, the other poison; in both of them one and the same person is concerned."

\end{quotebox}

\textbf{Why it matters.}

\begin{itemize}

\item Cicero narrows the case to Clodia-centered allegations.

\item He tries to strip away the public-crime frame and recast the prosecution as a private vendetta.

\end{itemize}

\textbf{From \texttt{Cicero In Defense of Marcus Caelius.pdf}, on Clodia's narrative.}

\begin{quotebox}\itshape

"This whole tale, for instance, belongs to an aging lady poet who has written a lot of comedies - but it has no plot and no ending."

\end{quotebox}

\textbf{Why it matters.}

\begin{itemize}

\item Cicero accuses the prosecution of theatrical fabrication while doing theatrical fabrication himself.

\item This is a perfect example of rhetorical projection.

\end{itemize}

\textbf{From \texttt{Cicero In Defense of Marcus Caelius.pdf}, section 70.}

\begin{quotebox}\itshape

"Members of the jury, my case is said and done ... This is the law under which this young man Caelius is being prosecuted not to exact punishment on behalf of the commonwealth but on behalf of the sexual perversions of this woman."

\end{quotebox}

\textbf{Why it matters.}

\begin{itemize}

\item Cicero returns to the seriousness of the \texttt{vis} law.

\item He argues that the state is being misused for private erotic revenge.

\item He fuses legal argument, misogyny, and comic performance.

\end{itemize}

\textbf{From the Week 12 lecture script, quoting Cicero's description of Catiline.}

\begin{quotebox}\itshape

"he could alter his character and direct it for the occasion, twisting and bending it in all directions"

\end{quotebox}

\textbf{Why it matters.}

\begin{itemize}

\item Cicero uses Catiline's adaptability to excuse Caelius' youthful association with him.

\item This shows how flexible courtroom rhetoric can be: last week's monster becomes this week's useful comparison.

\end{itemize}

\subsection{What the lecture wants you to notice}

\begin{calloutbox}[Key Takeaway]

\begin{itemize}

\item Cicero makes a public-crime case look like sex comedy.

\item He uses family, class, sexuality, and gender performance to assign criminality.

\item He relies on what looks very much like positivist thinking: good families produce good citizens, corrupt households produce corrupt actors.

\item The case also reveals a larger Roman shift toward treating unauthorized violence as violence against the state.

\end{itemize}

\end{calloutbox}

\subsection{Riggsby's contribution}

\textbf{The lecture uses Andrew Riggsby to argue.}

\begin{itemize}

\item Cicero replaces a judicial frame with a theatrical one.

\item Metaphor changes what outcomes feel "natural."

\item The law on \texttt{vis} helps show the Roman state's growing monopoly on legitimate force.

\end{itemize}

% START WEEK 12 MIDTERM FRAME
\subsection{Lecture frame in plain English}

The background topic shows that the Caelius trial is entangled with Egyptian diplomacy, elite rivalry, and the afterlife of the Catilinarian moment. The strategy topic argues that Cicero wins by theatrical substitution: he moves the case from state-danger narrative to comic-performance narrative. The criminality topic argues that Pro Caelio is about character, family, sexuality, and state control of violence at least as much as it is about factual guilt.

The charge does not make sense if you strip it out of late Republican politics and the Dio affair. It is the clearest Roman demonstration that genre and persona can do legal work. It brings together rhetoric, public/private wrong, gendered blame, and the Roman move toward stronger public claims over force. The public-private tension is part of the case, not background noise. The techniques change, but the struggle over interpretive frame remains familiar. The difference is partly one of overtness, not the disappearance of character-based reasoning.

\subsection{Key concepts in paragraph form}

\textbf{vis.} A public-law category of force or violence threatening civic order. It lets Rome frame certain unauthorized actions as attacks on the commonwealth. Cicero's defense depends on making this category look absurdly misapplied.

\textbf{Megalensia.} The festival setting during which public games and dramatic performances were happening. It gives Cicero a stage-like frame for turning the courtroom into a theater. The holiday setting is part of the speech's persuasive strategy, not a random detail.

\textbf{Theatrical substitution.} Replacing one interpretive frame with another more emotionally persuasive one. It makes jurors feel that one narrative genre fits the case more naturally than another. Cicero uses it to turn a public-law case into a comedy of youthful indiscretion and malicious plotting.

\textbf{Character evidence.} Reasoning from the kind of person someone is said to be rather than only from the specific act alleged. It makes social identity do evidentiary work. Cicero's portrait of Caelius and Clodia drives the entire defense.

\textbf{State monopoly on violence.} The gradual claim that legitimate force belongs increasingly to public authority. It turns private violence into a problem for the commonwealth. The law on vis is one of the clearest Roman markers of that shift.

\begin{trapbox}[MCQ Trap]
Do not describe Pro Caelio as only a sex scandal; the legal charge is public violence, even if Cicero works hard to privatize it. Do not say Cicero proves the facts first and performs second; the performance is one of the main proofs. Do not treat Clodia only as a historical figure; in the speech she is also a rhetorical construction designed to absorb blame.
\end{trapbox}
% END WEEK 12 MIDTERM FRAME

\subsection{Study questions answered}

\begin{enumerate}

\item \textbf{What are the Five Things?} The exam frame is Cicero, Pro Caelio, delivered in 56 BCE at Rome in Latin. Those facts matter because the speech belongs to the late Republic, to a law court on vis, and to a political culture in which rhetoric, festival atmosphere, and elite scandal all overlap.

\item \textbf{What are the background issues of the trial?} The lecture places the case inside Egyptian diplomacy, the Alexandrian embassy, the death of Dio, and the aftershocks of late Republican instability. A good answer makes clear that the trial is not merely a private quarrel but a public prosecution entangled with politics and foreign affairs.

\item \textbf{Who is on the prosecution team?} Atratinus, Publius Clodius, and Herennius Balbus form the prosecution side. The important point is not only the names but the fact that Cicero turns their case into something that looks dependent on Clodia's scandalized narrative.

\item \textbf{Who is on the defense team?} Caelius, Crassus, and Cicero appear for the defense, with Cicero serving as the climactic speaker. On the exam, connect that role to performance: Cicero is not simply arguing doctrine but controlling the whole emotional frame of the trial.

\item \textbf{Who is Cicero in this trial?} He is the final defense speaker and the architect of the speech's theatrical strategy. The lecture treats him as someone who wins not mainly by technical refutation but by making one story feel comic, natural, and socially legible while making the rival story feel ridiculous.

\item \textbf{How does Cicero characterize Caelius?} Cicero presents Caelius as talented, socially acceptable, and guilty only of ordinary youthful excess. That characterization matters because it turns character itself into evidence and makes the jurors feel that Caelius belongs inside elite normality rather than inside the category of dangerous offender.

\item \textbf{Who is Clodia and what role does she play?} Historically she is an elite Roman woman and Caelius' former lover; rhetorically she becomes the speech's central scapegoat. The lecture stresses that you must treat her as both a person and a constructed comic type, because Cicero wins partly by collapsing those two levels together.

\item \textbf{What is the relationship between Caelius and Catiline?} Caelius had connections to Catiline's circle, and that association is dangerous after the events of Week 11. Cicero's answer is to present Catiline as so deceptive and protean that even association with him no longer proves Caelius' guilt.

\item \textbf{What is Cicero's legal strategy and what theory does it resemble?} His strategy is to replace a strict public-law frame with character, genre, ridicule, and theatrical substitution. That resembles positivist or stereotype-based thinking because it asks the jury to decide what kind of people Caelius and Clodia are before deciding what happened.

\item \textbf{How are elements of this trial still visible today?} The lecture points to media spectacle, celebrity scandal, sexual double standards, character assassination, and metaphor-driven politics. The strongest comparison is that modern trials still compete in public culture through branding, persona, and story, even when formal doctrine pretends to be more controlled.

\end{enumerate}

\subsection{Deep study layer}

\begin{calloutbox}[What this lecture is arguing]
Week 12 argues that \texttt{Pro Caelio} is a criminal-defense speech in which genre does as much legal work as doctrine. Cicero does not mainly win by disproving every allegation in a technical way. He wins by relocating the case from the severe world of public violence into the comic world of sex scandal, youthful error, and malicious ex-lover narrative.
\end{calloutbox}

\begin{conceptbox}[Concept Bubble: \texttt{vis} / Megalensia / theatrical substitution]
\textbf{\texttt{vis}} marks certain force as a threat to public order and the commonwealth. \textbf{The Megalensia} gives Cicero a ready-made theatrical frame, since the jurors would otherwise be at public performances. \textbf{Theatrical substitution} replaces one interpretive frame with another: not armed public danger but comic intrigue; not state crisis but elite domestic scandal.
\end{conceptbox}

\begin{contextbox}[Passage Context: the main \texttt{Pro Caelio} quotations]
\texttt{[1]} is Cicero's opening to jurors in person during a holiday trial under a law designed for serious civic violence; its job is to inflate the legal frame only so he can ridicule its application to Caelius. The ``two charges'' passage narrows the case and privatizes it around Clodia. The ``aging lady poet'' passage converts Clodia into a comic type so that her credibility collapses with her genre. The \texttt{[70]} passage returns to the gravity of the \texttt{vis} statute and claims that a public law is being hijacked for private erotic revenge.
\end{contextbox}

\begin{comparisonbox}[Ancient vs Modern]
The strongest modern parallel is the media-saturated trial in which evidence competes with persona, scandal, sexual morality, and narrative branding. Modern courts constrain overt attacks on a woman's sexual character more than Cicero's Rome did, but public discourse still repeatedly turns legal controversies into stories about lifestyle, charisma, and ``who these people really are.''
\end{comparisonbox}

\begin{trapbox}[Exam Trap]
Do not say that Cicero simply proves Caelius innocent. The brilliance of the speech is that it often avoids fact-heavy proof altogether and instead wins by making one genre of explanation feel more natural than another.
\end{trapbox}

\section{Cross-Course Comparison Table}

\begin{longtable}{>{\raggedright\arraybackslash}p{0.15\textwidth}>{\raggedright\arraybackslash}p{0.21\textwidth}>{\raggedright\arraybackslash}p{0.23\textwidth}>{\raggedright\arraybackslash}p{0.26\textwidth}}
\toprule
Topic & Greek material & Roman material & Why it matters \\
\midrule
\endfirsthead
\toprule
Topic & Greek material & Roman material & Why it matters \\
\midrule
\endhead
Public/private distinction & Homicide can remain private & Many violent matters remain outside strong public prosecution for a long time & Modern criminal law is not universal \\
Self-help & Draco regulates it & Patron-client structure and weak policing preserve it & Ancient states do not fully monopolize force \\
Rhetoric & Lysias and Gorgias & Sallust and Cicero & Persuasion constructs criminality \\
Punishment & exile, blood vengeance, civic adjudication & exile, fines, execution, later stronger criminal courts & Prison is not the default ancient answer \\
Character versus act & seducer, adulterer, tyrant & Catiline, Clodia, Caelius & The person often matters more than the legal category \\
Insider versus outsider & household polluter, tyrant, seducer & citizen conspirator, aristocrat, enemy combatant & The threat often comes from within \\
Evidence and proof & witnesses, oaths, narrated household scenes & statutes, speeches, senatorial debates, reputation & Facts rarely appear without rhetorical framing \\
Emergency power & divine and kin pressure before stable institutions & \texttt{SCU}, \texttt{hostis}, courts on \texttt{vis} & Security crises put pressure on ordinary procedure \\
Gender and crime & Clytemnestra, adulterous wife, seduced household & Clodia, elite women, comic and sexualized blame & Gender often becomes a shortcut for explaining danger \\
Monopoly on violence & incomplete, contested, emerging & increasingly centralized in public law & A major course arc is the slow move from regulated self-help to stronger state claims \\
\bottomrule
\end{longtable}

\section{High-Yield Terms to Memorize}

\subsection{Greek}

\begin{itemize}

\item \texttt{nomos}

\item \texttt{dike}

\item \texttt{graphe}

\item lethal self-help

\item Areopagus

\item Socrates

\item rhetoric

\item \texttt{physis}

\item \texttt{nomos} as convention

\end{itemize}

\subsection{Roman}

\begin{itemize}

\item \texttt{imperium}

\item \texttt{intercessio}

\item \texttt{provocatio}

\item Twelve Tables

\item \texttt{praetor urbanus}

\item \texttt{praetor peregrinus}

\item \texttt{ius honorarium}

\item \texttt{senatus consultum}

\item \texttt{senatus consultum ultimum}

\item \texttt{quaestio}

\item \texttt{quaestio perpetua}

\item \texttt{repetundae}

\item \texttt{maiestas}

\item \texttt{ambitus}

\item \texttt{vis}

\item \texttt{hostis}

\item \texttt{ius respondendi}

\item \texttt{Corpus Iuris Civilis}

\item \texttt{honestiores}

\item \texttt{humiliores}

\item delict

\end{itemize}

% START GENERATED GLOSSARY
\section{Full Glossary Appendices}

\begin{calloutbox}[Glossary Note]
These glossary definitions are condensed study shorthand adapted from \texttt{course-materials/misc/Greek\_Law\_Wiki-45763.html} and \texttt{course-materials/misc/roman glossary.pdf}. Use them for quick review, not as a replacement for official course definitions.
\end{calloutbox}

\subsection{Greek Glossary}

\begin{itemize}

\item \textbf{Agon}: Formal dispute scene in Greek tragedy or comedy.

\item \textbf{Anakrisis}: Pre-trial hearing where the magistrate read the complaint and asked for admission or denial.

\item \textbf{Apagoge}: The procedure by which the volunteer prosecutor arrests the defendant first and then brings him to the competent official.

\item \textbf{Apagoge, ephegesis, endeixis, apographe, eisangelia, probole, dokimasia, euthynai}: Types of public cases that volunteer prosecutors could bring, often called "summary procedures.".

\item \textbf{Apographe}: The procedure by which the volunteer prosecutor lists property that is wrongfully held by the defendant, property that rightfully belongs to the State.

\item \textbf{Apotumpansimos}: The act of exposing the convicted individual by securing his neck, wrists, and ankles to a vertical wooden plank (another form of the death penalty).

\item \textbf{Archon}: Chief magistrate at Athens.

\item \textbf{Archon Basileus}: Archon responsible for religion, homicide, and deliberate wounding cases.

\item \textbf{Archon Eponymous}: Archon responsible for property, family matters, and key festivals; the year was named after him.

\item \textbf{Archon Polemarchus}: Archon responsible for cases involving non-Athenians.

\item \textbf{Areopagus}: Aristocratic council and court for intentional homicide and other grave offenses.

\item \textbf{Barathon}: The pit into which the convicted individual would be hurled (for death penalties).

\item \textbf{Boedromion \& Mounikhion}: Certain actions/cases could only be brought out between specific months.

\item \textbf{Boule}: Council of 500 citizens that prepared business for the assembly.

\item \textbf{De novo}: A Latin expression meaning "from scratch", "from the beginning".

\item \textbf{Delphinion}: Homicide cases were tried in this court when the defendant was asserting an affirmative defense (arguing that the homicide was excused on a grounds of assumption of risk, self-defense, defense of property, etc.).

\item \textbf{Diamartyria}: Literally "on the account of a testimony", "through testimony".

\item \textbf{Dikai emporikai}: Mercantile cases that applied to both citizens and non-citizens and could only be initiated if the dispute at issue pertained to an alleged breach of a written contract that either: had been concluded in the Athenian...

\item \textbf{Dikastai/dikastes}: Jurors.

\item \textbf{Dikastic oath}: An oath that each dikastes (juror) was required to swear in order to be eligible to serve as a juror.

\item \textbf{Dike}: Private suit or case brought by the wronged individual.

\item \textbf{Dike aikeias}: A charge of assault or battery (the Athenians didn't distinguish between the two).

\item \textbf{Dike biaien}: A cause of action for theft by force.

\item \textbf{Dike blabes}: A case for damage to property or breach of contract; a cause of action used to rectify/correct a variety of civil wrongs, e.g.

\item \textbf{Dike exoules}: An additional fine paid by the defendant to the state if the defendant failed to pay the plaintiff what was due.

\item \textbf{Dike heirgmou}: A charge for false imprisonment/confinement.

\item \textbf{Dike kakegorias}: An action brought out for slander/defamation.

\item \textbf{Dike klopes}: A cause of action for simple theft.

\item \textbf{Dokimasia}: A hearing where a candidate might be qualified from citizenship, public office or speaking in the Ekklesia.

\item \textbf{Eisagogeis}: Monetary cases related to loans and banking.

\item \textbf{Eisangelia}: Public denunciation procedure used for treason, betrayal, and other injuries to the state.

\item \textbf{Ek pronoias}: The Greek term meaning "intentionally," "with forethought," "with design," "purposely." A homicide that could be described by this term was punishable by death and the confiscation of the murderer's property.

\item \textbf{Ekklesia}: Assembly of Athenian citizens.

\item \textbf{Endeixis}: The procedure by which the volunteer prosecutor first explains the charge to the magistrate and is then authorized to make the arrest.

\item \textbf{Ephegesis}: The procedure by which the volunteer prosecutor leads the magistrate to the defendant for arrest.

\item \textbf{Ephetai}: A special court of 51 men over the age of 50 who tried homicide cases in four distinct courts: the Delphinion; the Palladion; the Prytaneion; the Phreatto.

\item \textbf{Epimeletai tou emporiou}: Cases involving the sale of grain in the whole sale market that was conducted in the Athenian port, the Piraeus, at the Emporion.

\item \textbf{Euthynai}: Review of performance in public office.

\item \textbf{Exegetai}: (="expounders") Religious officials who memorized sacred laws so that they could advise others about the laws.

\item \textbf{Graphe}: Public action brought on behalf of the community.

\item \textbf{Graphe moicheias}: A private case against a rapist.

\item \textbf{Graphe paranomon}: A "prosecution for illegalities" used to block decrees that contradicted or conflicted with established laws.

\item \textbf{Heken}: This term also means "intentionally." A homicide that could be described by this term was punishable by death and the confiscation of the murderer's property.

\item \textbf{Hekousios}: This term also means "intentionally," or "voluntarily." A homicide that could be described by this term was punishable by death and the confiscation of the murderer's property.

\item \textbf{Heliaia}: Judicial sitting of the Athenian assembly or people's court.

\item \textbf{Hierosylia}: The act of robbing from a temple (which could be punished by death).

\item \textbf{Hubris}: Abusive, arrogant, degrading conduct that could ground prosecution.

\item \textbf{Kakourgoi}: A cause of action for serious theft, e.g. thieves who stole at night, stole from a gymnasium, stole more than 10 drachmas from a harbor or stole more than 50 drachmas from anywhere else.

\item \textbf{Klepsydra}: A water clock that kept track of the time allotted for each speech.

\item \textbf{Kleroterion}: An allotment machine used to sort black and white balls for jury selection.

\item \textbf{Kyrios}: Legally responsible adult male guardian of a woman or child.

\item \textbf{Lex}: The Latin term for law or statute.

\item \textbf{Lex talionis}: (Latin) Law of retribution; an eye for an eye; follow for follow or rather, retaliatory retribution.

\item \textbf{Logographoi}: Speechwriters who composed court speeches for litigants.

\item \textbf{Martyria}: A testimony.

\item \textbf{Miasma}: Pollution caused by homicide that could spread through the community.

\item \textbf{Nautodikai}: Judge of sailors who handled cases involving Athenians who either lived overseas or sailed as mariners or merchants.

\item \textbf{Nomos/Nomoi}: Greek law or custom; later also written statute.

\item \textbf{Nomothetai}: Lawgivers or legislators who approved laws after assembly votes.

\item \textbf{Palladion}: Court used for unintentional homicides or when the victim was either a slave or foreigner.

\item \textbf{Paragraphe}: A procedural challenge; "prosecution in opposition", "counter prosecution"; a technical legal mechanism employed by defendants.

\item \textbf{Phreatto}: Where cases were heard when the accused had already been convicted and exiled for a prior homicide.

\item \textbf{Pinakion}: A wooden or bronze "ticket" that had each juror's name on it and one of the first ten letters of the alphabet.

\item \textbf{Poine}: A payment.

\item \textbf{Polis}: Greek city-state understood as a political community, not just a place.

\item \textbf{Probole}: A preliminary hearing at the Ekklesia regarding official misconduct.

\item \textbf{Prodikasiai}: Pre-trial conferences; a procedural stage for homicide trials.

\item \textbf{Prodosia}: A public-wrong charge used against generals for losing ships, troops, or territory.

\item \textbf{Prosklesis}: A summons or "call to court"; the summons had to alert the defendant to three facts: 1.

\item \textbf{Prytaneia}: A filing fee which the plaintiff/prosecutor had to pay.

\item \textbf{Prytaneion}: Court used for cases where death was caused by an unknown person (i.e.

\item \textbf{Pseudomartyrion}: A false testimony or perjury. A dike pseudomartyrion is a charge/prosecution for bearing false testimony.

\item \textbf{Strategoi}: Board of generals with military jurisdiction.

\item \textbf{Sychophants}: Volunteer prosecutors who brought cases for profit, prestige, or leverage.

\item \textbf{Synegroi}: (="supporting speakers") Public prosecutors who brought actions against men who were acting in an official capacity.

\item \textbf{The Eleven}: Dealt with suits involving kakourgoi and cases that required the defendant to be incarcerated while awaiting trial.

\item \textbf{The Forty}: Judges who heard majority of the private cases brought by means of a dike (except under the jurisdiction of the Archon, Thesmothetai, or some other specialized court).

\item \textbf{Thesmothetai}: Officials responsible for many graphe and dike cases.

\item \textbf{Trauma ek pronoias}: Wounding with intent to kill; deliberate armed battery.

\item \textbf{Xenodikai}: Judges of foreigners (Abolished around 350 BCE).

\end{itemize}

\subsection{Roman Glossary}

\begin{itemize}

\item \textbf{Abigeatus}: Rustling; a specific subspecies of theft.

\item \textbf{Absolvo}: Acquittal.

\item \textbf{Actio de posito et suspenso}: "an action relating to something that has been placed or hung out".

\item \textbf{Actio de rebus effusis vel deiectis}: "an action relating to things which have been thrown out" This action was carried out against the owner of the property from which things were thrown or poured onto a street.

\item \textbf{Actio servi corrupti}: This action concerned situations where a person influenced another's slave in a negative manner.

\item \textbf{Actus reus}: Criminal act; an essential element for a crime in Roman law.

\item \textbf{Ad agendum}: To prepare a case for the court (Latin meaning: agere is the verb used to plead a case before the court).

\item \textbf{Ad cavendum}: To draft documents; the process of drafting written formulae for lawsuits or business transactions.

\item \textbf{Ad respondendum}: To answer legal questions and give legal advice; the private function of a lawyer to explanation to a praetor, aedile, judge, or layperson what a particular law meant.

\item \textbf{Aedile}: Roman magistrate who supervised archives, public works, streets, and the marketplace.

\item \textbf{Ambitus}: Criminal conduct related to elections such as bribery and other types of electoral corruption.

\item \textbf{Apud iudicem}: The time a case would be heard by a judge.

\item \textbf{Augurs}: Augurs were in charge of the auspices and augury. They had the power to delay or prohibit public business.

\item \textbf{Calumnia}: Bringing false criminal charges as an accuser.

\item \textbf{Censor}: Magistrate who ran the census and assessed the social and moral standing of Roman citizens.

\item \textbf{Centumviri (100 men)}: Large juries used for disputes involving inheritance of large estates of nobles or the rich.

\item \textbf{Civis}: Latin for citizen.

\item \textbf{Cognitio extraordinario / Cognitio extra ordinem}: Extraordinary post-classical procedure in which one official heard the whole case instead of splitting it into separate phases.

\item \textbf{Comitia}: An assembly.

\item \textbf{Comitia centuriata}: The assembly compromising all Roman citizens organized into "centuries". It was organized by wealth.

\item \textbf{Comitia curiata}: During the Republic it served mostly to witness and perhaps authorize wills and adoptions. It operated in the Monarchy but we don't know the details.

\item \textbf{Comitia tributa}: The comitia tributa was organized by the location of residence and they vote in an order determined by lot.

\item \textbf{Concilium plebis}: An assembly of all plebeians organized by tribe and presided over by the tribunes of the plebs.

\item \textbf{Condemnatio}: This is another part of the typical formula. It is the instruction to judge to either find the defendant liable or not based upon an application of the formula to his findings of fact.

\item \textbf{Condemno}: A conviction.

\item \textbf{Contiones}: Public meetings Where debating and public discussion took place.

\item \textbf{Cursus honorum}: A regular order of elected offices which was established after the Second Punic War.

\item \textbf{Damnum iniuria datum}: Property damage.

\item \textbf{De residuis}: Embezzlement of entrusted money.

\item \textbf{Decreta}: Legal decisions rendered by the emperor in any given case that had been brought before him.

\item \textbf{Delictum}: Crime or unintentional tort.

\item \textbf{Dolus}: Criminal intent.

\item \textbf{Edictum perpetuum}: Julianus, the jurist under the reign of Emperor Hadrian (117-138 CE) put the praetorian and aedilician edicts into their final forms. Later classical Roman jurists referred to Julianus' frozen version of the...

\item \textbf{Effractores}: Burglars who broke into apartments.

\item \textbf{Equites (singular: eques )}: They bid on contracts with the senate or with magistrates to carryout public business like the collection of taxes (publicani).

\item \textbf{Exceptio}: An exceptio is an example of one of the legal institutions that Roman praetors created. It is a method of protecting rights kind of like a defense to an action.

\item \textbf{Expilatores}: Thieves who ransacked homes in the countryside.

\item \textbf{Falsum}: Forgery.

\item \textbf{Flagrante delicto}: (literally,"in blazing offence"); caught red handed.

\item \textbf{Furtum manifestum}: The thief is caught in the action. The penalty was to be beaten and then the perpetrator was handed over to the victim of the theft as a slave or bondsman.

\item \textbf{Furtum nec manifestum}: The thief is not caught in the action.

\item \textbf{Glossa}: Comments and laws.

\item \textbf{Honestiores}: More privileged citizens (for example army veterans).

\item \textbf{Imperium}: The power to command held by senior Roman magistrates.

\item \textbf{In iure}: The first phase of Roman procedure, conducted before the praetor.

\item \textbf{Iniuria}: Personal injury to reputation, dignity, honor, or bodily integrity.

\item \textbf{Inscriptio}: The court president wrote down an inscriptio for the prosecutor to sign.

\item \textbf{Institutiones}: An official textbook on Roman law and a functional code that could be consulted for subsequent decisions.

\item \textbf{Intentio}: This is another part of the typical formula (written instruction). It is the statement of the plaintiff's claim and the most important section of every formula.

\item \textbf{Intercessio}: Veto power used against magistrates of equal or lower status.

\item \textbf{Interdictio aquae et ignis}: Banishment/exile; literally "prohibition of water and fire".

\item \textbf{Iudex}: Latin for "judge". One of the most important officials involved in the development of the legis actiones system.

\item \textbf{Iudex qui litem suam fecit}: "a judge who has made a case his own". One of the quasi-delicts.

\item \textbf{Iudicis nominatio}: This is one part of the typical formula which nominates a judge.

\item \textbf{Ius Aelianum}: A collection of judicial formulae for lawsuits, dating to ca. 200 BCE.

\item \textbf{Ius Flavianum}: A collection of formulae for lawsuits, dating to ca. 300 BCE.

\item \textbf{Ius gentium}: Literally, "law of the peoples/nations". The praetor peregrinus used this law between a Roman citizen and a foreigner or between two foreigners who were involved in a dispute on Roman soil.

\item \textbf{Ius honorarium}: The law that resulted from the praetors' edicts. It was used to distinguish law based on the praetor's edict from more formal legislation.

\item \textbf{Ius occidendi}: The right of killing; stems from the power of the paterfamilias over members of his family.

\item \textbf{Ius privatum}: Included rules of property, succession, contracts, and laws relating to the family; deals with the interests of separate persons (according to the Roman Jurist Ulpian).

\item \textbf{Ius publicum}: A body of laws that deals with interests of the entire community (according to Roman Jurist Ulpian).

\item \textbf{Ius respondendi ex auctoritate principis}: (sometimes referred to simply as the ius respondendi ); the law of responding from the authority of the princeps (emperor).

\item \textbf{Ius vocatio}: The act of bringing a defendant before the praetor by a formal summons initiated in the in iure phase.

\item \textbf{Laudatores}: Character witnesses.

\item \textbf{Leges}: Statutes adopted by the populus romanus; designed to honour the honour the mos maiorum.

\item \textbf{Legis actiones}: The early form of civil procedure; "suits/actions of law".

\item \textbf{Lenocinium}: Any conduct that facilitated sexual crime.

\item \textbf{Lex Aquilia}: Plebiscite governing wrongful damage to property, especially injury to another person's slaves or animals.

\item \textbf{Lex Calpurnia}: Law of 149 BCE that created the first permanent jury court and helped launch the quaestiones perpetuae.

\item \textbf{Lex Canuleia}: A law passed in 445 BCE which gave plebeians the right of intermarriage with the patricians.

\item \textbf{Lex Cornelia de iniuriis}: A law that created what amounted to a criminal cause of action against persons for hittng others or for forcefully breaking into another's home.

\item \textbf{Lex Cornelia nummaria}: The earliest Roman law to criminalize the counterfeiting of money.

\item \textbf{Lex Fabia}: A law that criminalized the sale of a Roman citizen, freedman, or slave who belonged to someone else.

\item \textbf{Lex Hortensia}: Law of 287 BCE making plebiscites binding on all Roman citizens.

\item \textbf{Lex Iulia de adulteriis coercendis}: Augustan law punishing adultery and giving the paterfamilias limited rights in flagrante delicto.

\item \textbf{Lex Ogulnia}: (300 BCE) The law that allowed plebeians to be pontifices.

\item \textbf{Lex Pompeia}: (55 BCE) The first statute that criminalized parricide as an offense distinct from ordinary murder.

\item \textbf{Lex talionis}: A law on retribution; essentially the equivalent of "an eye for an eye".

\item \textbf{Licinian Sextian law}: (367 BCE) this law allowed one of the two consuls elected every year had to be a plebeian.

\item \textbf{Litis contestatio}: The appearance of parties before the praetor to initiate a suit.

\item \textbf{Maiestas}: Treason--Maiestas included a number of illegal activities: armed assault on the State; gathering soldiers or conducting warfare without the State's permission; desertion.

\item \textbf{Maius imperium}: The "greater power to command", held by a more senior magistrate (i.e.

\item \textbf{Mancipatio}: The act of transferring res mancipi e.g. land, slaves, cattle, horses, mules, donkeys.

\item \textbf{Mens rea}: "a guilty mind".

\item \textbf{Mos maiorum}: Customs of their (Romans') ancestors.

\item \textbf{Nominis delatio}: =(denunciation of the name) The trial began with a nominis delatio attended by both prosecutor and defendant.

\item \textbf{Novella}: New laws that were adopted after the Corpus Iuris had been written.

\item \textbf{Obligatio/obligationes}: This entailed some relationship between two persons in which one was a creditor and the other debtor.

\item \textbf{Paterfamilias}: The father of the family; the head of the household.

\item \textbf{Peculatus}: A kind of embezzlement; any wrongful appropriation, diversion, or conversion of sacred, religious, or public funds.

\item \textbf{Per formulam}: The formulary procedure which replaced the legis actiones in the beginning of the 2nd century B.C.

\item \textbf{Perduellio}: This term also refers to treasonous conduct but it was used more in the earlier Republic.

\item \textbf{Plagium}: Kidnapping.

\item \textbf{Plebiscita (singular: plebiscitum )}: Enactments of the concilium plebis that became binding on all Rome after the lex Hortensia.

\item \textbf{Pontifices}: Pontiffs (priests) had access to the archives containing the specialized legal forms and phrases that were necessary to conduct a law suit.

\item \textbf{Populus Romanus}: The Roman people.

\item \textbf{Praefectus vigilum}: Prefect of the Night Watch.

\item \textbf{Praescriptio}: A general comment concerning the wrong.

\item \textbf{Praetor}: Roman magistrate responsible for the judiciary, annual edicts, and legal formulae.

\item \textbf{Praetor peregrinus}: Praetor who handled disputes involving foreigners.

\item \textbf{Praetor urbanus}: Praetor who handled disputes involving Roman citizens.

\item \textbf{Praevaricatio}: A criminal offence for an accuser to hide legitimate charges.

\item \textbf{Provocatio}: Right to appeal a magistrate's summary judgment to the Roman people.

\item \textbf{Quaestio de repetundiis}: Court for extortion by provincial governors.

\item \textbf{Quaestiones perpetuae}: Permanent Roman criminal courts for recurring categories of public wrong.

\item \textbf{Quaestors}: Financial officials.

\item \textbf{Rapina}: Robbery with violence.

\item \textbf{Recuperatores}: Recoverers three to five judges instead of a single iudex.

\item \textbf{Repetundae}: Provincial extortion and abuse of office for private gain.

\item \textbf{Restitutio}: One of the legal institutions created by the Roman praetors.

\item \textbf{Saccularii}: Thieves who slit pockets and purses to empty their contents.

\item \textbf{Sanctio}: The negative consequences which resulted from a person's violation of the praescriptio.

\item \textbf{Sciens dolo malo}: "knowingly and with malicious intent".

\item \textbf{Senatus Consulta}: Advice or resolutions issued by the Senate.

\item \textbf{Sententia}: The opinion of the judge ( iudex ).

\item \textbf{Stellionatus}: Any conduct that was dishonest or fraudulent and did not fit neatly into another defined criminal category could be dealt as stellionatus.

\item \textbf{Stipulatio}: A face to face contract that originally required one party to ask "Do you promise" and the other to respond "I promise".

\item \textbf{Stuprum}: An illegal crime which occurred when a man acting sciens dolo malo had sexual relations with a "respectable" unmarried girl or woman or boy.

\item \textbf{Tergiversatio}: A criminal offence for an accuser to abandon charges once formally begun.

\item \textbf{Tresviri capitalis}: Handled cases involving theft.

\item \textbf{Tripertitia}: A major work published by Sextus Aelius Paetus. It was the first attempt at a systematic treatment of Roman law.

\item \textbf{Vis}: Force or violence treated as a public crime, often punished by exile or death.

\end{itemize}
% END GENERATED GLOSSARY

\section{Final Exam Focus: What to prioritize in the last review pass}

\begin{calloutbox}[Exam Callout]
Use this weighting for the last pass before the exam: master Roman legal structure first, then Catiline and \texttt{Pro Caelio}, then use the Greek units to define terms and build comparison answers.
\end{calloutbox}

\subsection{Highest priority}

\begin{calloutbox}[Key Takeaway]

\begin{itemize}

\item Week 10 Roman law terms, institutions, assemblies, procedures, and codification

\item Week 11 Catiline chronology, \texttt{SCU}, \texttt{hostis}, Caesar versus Cato

\item Week 12 \texttt{Pro Caelio}, \texttt{vis}, Cicero's theatrical strategy, Clodia, and character-based criminality

\end{itemize}

\end{calloutbox}

\subsection{Second priority}

\begin{calloutbox}[Key Takeaway]

\begin{itemize}

\item Week 7 Draco's law and Lysias' defense strategy

\item Week 8 Socrates versus Polus and Callicles

\item Week 6 criminology definitions and theories

\end{itemize}

\end{calloutbox}

\subsection{Foundation priority}

\begin{calloutbox}[Key Takeaway]

\begin{itemize}

\item Week 1 public/private distinction, exile versus prison, Foucault

\item Week 2 chronology

\item Week 3 Greek-law vocabulary

\item Week 4 \texttt{Oresteia} as movement from vengeance toward civic justice

\end{itemize}

\end{calloutbox}

\section{Final Checklist}

\begin{itemize}

\item Can you explain why murder is not simply a "crime" in the same way in classical Athens?

\item Can you explain the difference between a homicide exception and a death penalty?

\item Can you define \texttt{vis} and explain why it matters in both Week 11 and Week 12?

\item Can you explain why the Twelve Tables matter for Roman legal history?

\item Can you explain the difference between \texttt{leges}, \texttt{plebiscita}, praetorian edicts, and juristic writings?

\item Can you explain why Caesar and Cato disagree about the Catilinarian conspirators?

\item Can you explain how Cicero turns a criminal prosecution into comedy?

\item Can you connect ancient material to modern criminological theories without claiming the two worlds are identical?

\end{itemize}

If you can answer those clearly, you are covering the spine of the course.

% START GENERATED GLOSSARY DRILL
\section{Glossary Practice Questions}

\subsection{Greek Glossary Drill}

\begin{enumerate}

\item Given the definition, what is the correct greek term? "Archon responsible for property, family matters, and key festivals; the year was named after him."
\begin{itemize}
\item A. Archon Eponymous
\item B. Boule
\item C. Archon Polemarchus
\item D. Areopagus
\end{itemize}

\item Given the definition, what is the correct greek term? "Legally responsible adult male guardian of a woman or child."
\begin{itemize}
\item A. Logographoi
\item B. Kyrios
\item C. Nomos/Nomoi
\item D. Polis
\end{itemize}

\item Given the definition, what is the correct greek term? "Council of 500 citizens that prepared business for the assembly."
\begin{itemize}
\item A. Boule
\item B. Heliaia
\item C. Ekklesia
\item D. Nomothetai
\end{itemize}

\item Given the definition, what is the correct greek term? "Private suit or case brought by the wronged individual."
\begin{itemize}
\item A. Dike
\item B. Archon Basileus
\item C. Archon
\item D. Trauma ek pronoias
\end{itemize}

\item Given the definition, what is the correct greek term? "Lawgivers or legislators who approved laws after assembly votes."
\begin{itemize}
\item A. Strategoi
\item B. The Eleven
\item C. The Forty
\item D. Nomothetai
\end{itemize}

\item Given the definition, what is the correct greek term? "Board of generals with military jurisdiction."
\begin{itemize}
\item A. Thesmothetai
\item B. The Eleven
\item C. Strategoi
\item D. The Forty
\end{itemize}

\item Given the definition, what is the correct greek term? "Chief magistrate at Athens."
\begin{itemize}
\item A. Archon Basileus
\item B. Archon
\item C. Archon Polemarchus
\item D. Archon Eponymous
\end{itemize}

\item Given the definition, what is the correct greek term? "A procedural challenge; "prosecution in opposition", "counter prosecution"; a technical legal mechanism employed by defendants."
\begin{itemize}
\item A. Paragraphe
\item B. Graphe
\item C. Prosklesis
\item D. Eisangelia
\end{itemize}

\item Given the definition, what is the correct greek term? "Archon responsible for religion, homicide, and deliberate wounding cases."
\begin{itemize}
\item A. Archon Polemarchus
\item B. Archon Eponymous
\item C. Archon Basileus
\item D. Areopagus
\end{itemize}

\item Given the definition, what is the correct greek term? "Public action brought on behalf of the community."
\begin{itemize}
\item A. Trauma ek pronoias
\item B. Graphe
\item C. Archon
\item D. Dike
\end{itemize}

\item Given the definition, what is the correct greek term? "Archon responsible for cases involving non-Athenians."
\begin{itemize}
\item A. Areopagus
\item B. Archon Polemarchus
\item C. Boule
\item D. Ekklesia
\end{itemize}

\item Given the definition, what is the correct greek term? "A summons or "call to court"; the summons had to alert the defendant to three facts: 1."
\begin{itemize}
\item A. Eisangelia
\item B. Dike
\item C. Graphe
\item D. Prosklesis
\end{itemize}

\item Given the definition, what is the correct greek term? "Speechwriters who composed court speeches for litigants."
\begin{itemize}
\item A. Polis
\item B. Logographoi
\item C. Nomos/Nomoi
\item D. Hubris
\end{itemize}

\item Given the definition, what is the correct greek term? "Judges who heard majority of the private cases brought by means of a dike (except under the jurisdiction of the Archon, Thesmothetai, or some other specialized court)."
\begin{itemize}
\item A. Dikastai/dikastes
\item B. The Forty
\item C. Thesmothetai
\item D. Kyrios
\end{itemize}

\item Given the definition, what is the correct greek term? "Aristocratic council and court for intentional homicide and other grave offenses."
\begin{itemize}
\item A. Boule
\item B. Ekklesia
\item C. Heliaia
\item D. Areopagus
\end{itemize}

\item Given the definition, what is the correct greek term? "Judicial sitting of the Athenian assembly or people's court."
\begin{itemize}
\item A. The Eleven
\item B. Heliaia
\item C. Nomothetai
\item D. Strategoi
\end{itemize}

\item Given the definition, what is the correct greek term? "Types of public cases that volunteer prosecutors could bring, often called "summary procedures."."
\begin{itemize}
\item A. Anakrisis
\item B. Apagoge
\item C. Apographe
\item D. Apagoge, ephegesis, endeixis, apographe, eisangelia, probole, dokimasia, euthynai
\end{itemize}

\item Given the definition, what is the correct greek term? "Public denunciation procedure used for treason, betrayal, and other injuries to the state."
\begin{itemize}
\item A. Dike
\item B. Trauma ek pronoias
\item C. Eisangelia
\item D. Graphe
\end{itemize}

\item Given the definition, what is the correct greek term? "Greek city-state understood as a political community, not just a place."
\begin{itemize}
\item A. Hubris
\item B. Polis
\item C. Apagoge, ephegesis, endeixis, apographe, eisangelia, probole, dokimasia, euthynai
\item D. Miasma
\end{itemize}

\item Given the definition, what is the correct greek term? "Dealt with suits involving kakourgoi and cases that required the defendant to be incarcerated while awaiting trial."
\begin{itemize}
\item A. The Forty
\item B. Dikastai/dikastes
\item C. Thesmothetai
\item D. The Eleven
\end{itemize}

\item Given the definition, what is the correct greek term? "Greek law or custom; later also written statute."
\begin{itemize}
\item A. Polis
\item B. Hubris
\item C. Nomos/Nomoi
\item D. Miasma
\end{itemize}

\item Given the definition, what is the correct greek term? "Pre-trial hearing where the magistrate read the complaint and asked for admission or denial."
\begin{itemize}
\item A. Apagoge
\item B. Anakrisis
\item C. Paragraphe
\item D. Apographe
\end{itemize}

\item Given the definition, what is the correct greek term? "Assembly of Athenian citizens."
\begin{itemize}
\item A. Ekklesia
\item B. Heliaia
\item C. Nomothetai
\item D. Strategoi
\end{itemize}

\item Given the definition, what is the correct greek term? "Pollution caused by homicide that could spread through the community."
\begin{itemize}
\item A. Apagoge, ephegesis, endeixis, apographe, eisangelia, probole, dokimasia, euthynai
\item B. Apagoge
\item C. Miasma
\item D. Anakrisis
\end{itemize}

\item Given the definition, what is the correct greek term? "The procedure by which the volunteer prosecutor arrests the defendant first and then brings him to the competent official."
\begin{itemize}
\item A. Apagoge
\item B. Paragraphe
\item C. Apographe
\item D. Prosklesis
\end{itemize}

\item Given the definition, what is the correct greek term? "Jurors."
\begin{itemize}
\item A. Nomos/Nomoi
\item B. Kyrios
\item C. Logographoi
\item D. Dikastai/dikastes
\end{itemize}

\item Given the definition, what is the correct greek term? "Officials responsible for many graphe and dike cases."
\begin{itemize}
\item A. Kyrios
\item B. Dikastai/dikastes
\item C. Thesmothetai
\item D. Logographoi
\end{itemize}

\item Given the definition, what is the correct greek term? "Wounding with intent to kill; deliberate armed battery."
\begin{itemize}
\item A. Trauma ek pronoias
\item B. Archon Basileus
\item C. Archon
\item D. Archon Eponymous
\end{itemize}

\item Given the definition, what is the correct greek term? "Abusive, arrogant, degrading conduct that could ground prosecution."
\begin{itemize}
\item A. Hubris
\item B. Apagoge, ephegesis, endeixis, apographe, eisangelia, probole, dokimasia, euthynai
\item C. Anakrisis
\item D. Miasma
\end{itemize}

\item Given the definition, what is the correct greek term? "The procedure by which the volunteer prosecutor lists property that is wrongfully held by the defendant, property that rightfully belongs to the State."
\begin{itemize}
\item A. Apographe
\item B. Prosklesis
\item C. Paragraphe
\item D. Eisangelia
\end{itemize}

\end{enumerate}

\subsection{Roman Glossary Drill}

\begin{enumerate}

\item Given the definition, what is the correct roman term? "Roman magistrate responsible for the judiciary, annual edicts, and legal formulae."
\begin{itemize}
\item A. Praetor urbanus
\item B. Praetor
\item C. Praetor peregrinus
\item D. Quaestors
\end{itemize}

\item Given the definition, what is the correct roman term? "The assembly compromising all Roman citizens organized into "centuries". It was organized by wealth."
\begin{itemize}
\item A. Senatus Consulta
\item B. Plebiscita (singular: plebiscitum )
\item C. Comitia tributa
\item D. Comitia centuriata
\end{itemize}

\item Given the definition, what is the correct roman term? "Criminal act; an essential element for a crime in Roman law."
\begin{itemize}
\item A. Mens rea
\item B. Actus reus
\item C. Dolus
\item D. Absolvo
\end{itemize}

\item Given the definition, what is the correct roman term? "Criminal conduct related to elections such as bribery and other types of electoral corruption."
\begin{itemize}
\item A. Delictum
\item B. Ambitus
\item C. Vis
\item D. Stuprum
\end{itemize}

\item Given the definition, what is the correct roman term? "Statutes adopted by the populus romanus; designed to honour the honour the mos maiorum."
\begin{itemize}
\item A. Ius honorarium
\item B. Legis actiones
\item C. Leges
\item D. Ius gentium
\end{itemize}

\item Given the definition, what is the correct roman term? "An illegal crime which occurred when a man acting sciens dolo malo had sexual relations with a "respectable" unmarried girl or woman or boy."
\begin{itemize}
\item A. Delictum
\item B. Stuprum
\item C. Calumnia
\item D. Peculatus
\end{itemize}

\item Given the definition, what is the correct roman term? "Magistrate who ran the census and assessed the social and moral standing of Roman citizens."
\begin{itemize}
\item A. Praetor urbanus
\item B. Praetor peregrinus
\item C. Censor
\item D. Praetor
\end{itemize}

\item Given the definition, what is the correct roman term? "Customs of their (Romans') ancestors."
\begin{itemize}
\item A. Mos maiorum
\item B. Paterfamilias
\item C. Interdictio aquae et ignis
\item D. Repetundae
\end{itemize}

\item Given the definition, what is the correct roman term? "Veto power used against magistrates of equal or lower status."
\begin{itemize}
\item A. Praetor
\item B. Censor
\item C. Intercessio
\item D. Aedile
\end{itemize}

\item Given the definition, what is the correct roman term? "A conviction."
\begin{itemize}
\item A. Provocatio
\item B. Iudex
\item C. Condemno
\item D. In iure
\end{itemize}

\item Given the definition, what is the correct roman term? ""a guilty mind"."
\begin{itemize}
\item A. Mens rea
\item B. Condemno
\item C. Absolvo
\item D. Dolus
\end{itemize}

\item Given the definition, what is the correct roman term? "Force or violence treated as a public crime, often punished by exile or death."
\begin{itemize}
\item A. Peculatus
\item B. Delictum
\item C. Vis
\item D. Stuprum
\end{itemize}

\item Given the definition, what is the correct roman term? "The law that resulted from the praetors' edicts. It was used to distinguish law based on the praetor's edict from more formal legislation."
\begin{itemize}
\item A. Ius gentium
\item B. Ius honorarium
\item C. Quaestiones perpetuae
\item D. Cognitio extraordinario / Cognitio extra ordinem
\end{itemize}

\item Given the definition, what is the correct roman term? "Treason--Maiestas included a number of illegal activities: armed assault on the State; gathering soldiers or conducting warfare without the State's permission; desertion."
\begin{itemize}
\item A. Ambitus
\item B. Maiestas
\item C. Vis
\item D. Stuprum
\end{itemize}

\item Given the definition, what is the correct roman term? "Praetor who handled disputes involving Roman citizens."
\begin{itemize}
\item A. Quaestors
\item B. Praetor peregrinus
\item C. Concilium plebis
\item D. Praetor urbanus
\end{itemize}

\item Given the definition, what is the correct roman term? "Literally, "law of the peoples/nations". The praetor peregrinus used this law between a Roman citizen and a foreigner or between two foreigners who were involved in a dispute on Roman soil."
\begin{itemize}
\item A. Ius gentium
\item B. Lex Calpurnia
\item C. Quaestiones perpetuae
\item D. Cognitio extraordinario / Cognitio extra ordinem
\end{itemize}

\item Given the definition, what is the correct roman term? "Financial officials."
\begin{itemize}
\item A. Comitia tributa
\item B. Quaestors
\item C. Concilium plebis
\item D. Comitia centuriata
\end{itemize}

\item Given the definition, what is the correct roman term? "Roman magistrate who supervised archives, public works, streets, and the marketplace."
\begin{itemize}
\item A. Praetor urbanus
\item B. Praetor
\item C. Censor
\item D. Aedile
\end{itemize}

\item Given the definition, what is the correct roman term? "Bringing false criminal charges as an accuser."
\begin{itemize}
\item A. Lex Iulia de adulteriis coercendis
\item B. Calumnia
\item C. Mens rea
\item D. Actus reus
\end{itemize}

\item Given the definition, what is the correct roman term? "Advice or resolutions issued by the Senate."
\begin{itemize}
\item A. Leges
\item B. Senatus Consulta
\item C. Imperium
\item D. Legis actiones
\end{itemize}

\item Given the definition, what is the correct roman term? "Enactments of the concilium plebis that became binding on all Rome after the lex Hortensia."
\begin{itemize}
\item A. Senatus Consulta
\item B. Plebiscita (singular: plebiscitum )
\item C. Leges
\item D. Imperium
\end{itemize}

\item Given the definition, what is the correct roman term? "Plebiscite governing wrongful damage to property, especially injury to another person's slaves or animals."
\begin{itemize}
\item A. Interdictio aquae et ignis
\item B. Lex Aquilia
\item C. Paterfamilias
\item D. Mos maiorum
\end{itemize}

\item Given the definition, what is the correct roman term? "Provincial extortion and abuse of office for private gain."
\begin{itemize}
\item A. Vis
\item B. Repetundae
\item C. Ambitus
\item D. Maiestas
\end{itemize}

\item Given the definition, what is the correct roman term? "Law of 149 BCE that created the first permanent jury court and helped launch the quaestiones perpetuae."
\begin{itemize}
\item A. Mos maiorum
\item B. Lex Hortensia
\item C. Lex Aquilia
\item D. Lex Calpurnia
\end{itemize}

\item Given the definition, what is the correct roman term? "Law of 287 BCE making plebiscites binding on all Roman citizens."
\begin{itemize}
\item A. Paterfamilias
\item B. Mos maiorum
\item C. Lex Aquilia
\item D. Lex Hortensia
\end{itemize}

\item Given the definition, what is the correct roman term? "An assembly of all plebeians organized by tribe and presided over by the tribunes of the plebs."
\begin{itemize}
\item A. Comitia tributa
\item B. Plebiscita (singular: plebiscitum )
\item C. Comitia centuriata
\item D. Concilium plebis
\end{itemize}

\item Given the definition, what is the correct roman term? "Right to appeal a magistrate's summary judgment to the Roman people."
\begin{itemize}
\item A. Censor
\item B. Aedile
\item C. Provocatio
\item D. Intercessio
\end{itemize}

\item Given the definition, what is the correct roman term? "The father of the family; the head of the household."
\begin{itemize}
\item A. Paterfamilias
\item B. Maiestas
\item C. Interdictio aquae et ignis
\item D. Repetundae
\end{itemize}

\item Given the definition, what is the correct roman term? "Augustan law punishing adultery and giving the paterfamilias limited rights in flagrante delicto."
\begin{itemize}
\item A. Lex Iulia de adulteriis coercendis
\item B. Actus reus
\item C. Dolus
\item D. Mens rea
\end{itemize}

\item Given the definition, what is the correct roman term? "Acquittal."
\begin{itemize}
\item A. Absolvo
\item B. In iure
\item C. Iudex
\item D. Condemno
\end{itemize}

\item Given the definition, what is the correct roman term? "Praetor who handled disputes involving foreigners."
\begin{itemize}
\item A. Praetor peregrinus
\item B. Comitia centuriata
\item C. Concilium plebis
\item D. Quaestors
\end{itemize}

\item Given the definition, what is the correct roman term? "Permanent Roman criminal courts for recurring categories of public wrong."
\begin{itemize}
\item A. Lex Calpurnia
\item B. Lex Aquilia
\item C. Quaestiones perpetuae
\item D. Lex Hortensia
\end{itemize}

\item Given the definition, what is the correct roman term? "A kind of embezzlement; any wrongful appropriation, diversion, or conversion of sacred, religious, or public funds."
\begin{itemize}
\item A. Actus reus
\item B. Peculatus
\item C. Lex Iulia de adulteriis coercendis
\item D. Calumnia
\end{itemize}

\item Given the definition, what is the correct roman term? "Crime or unintentional tort."
\begin{itemize}
\item A. Delictum
\item B. Calumnia
\item C. Peculatus
\item D. Lex Iulia de adulteriis coercendis
\end{itemize}

\item Given the definition, what is the correct roman term? "The first phase of Roman procedure, conducted before the praetor."
\begin{itemize}
\item A. Provocatio
\item B. Intercessio
\item C. In iure
\item D. Iudex
\end{itemize}

\item Given the definition, what is the correct roman term? "Latin for "judge". One of the most important officials involved in the development of the legis actiones system."
\begin{itemize}
\item A. Provocatio
\item B. Intercessio
\item C. Iudex
\item D. Aedile
\end{itemize}

\item Given the definition, what is the correct roman term? "Extraordinary post-classical procedure in which one official heard the whole case instead of splitting it into separate phases."
\begin{itemize}
\item A. Cognitio extraordinario / Cognitio extra ordinem
\item B. Lex Hortensia
\item C. Lex Calpurnia
\item D. Quaestiones perpetuae
\end{itemize}

\item Given the definition, what is the correct roman term? "The early form of civil procedure; "suits/actions of law"."
\begin{itemize}
\item A. Cognitio extraordinario / Cognitio extra ordinem
\item B. Legis actiones
\item C. Ius gentium
\item D. Ius honorarium
\end{itemize}

\item Given the definition, what is the correct roman term? "Banishment/exile; literally "prohibition of water and fire"."
\begin{itemize}
\item A. Ambitus
\item B. Interdictio aquae et ignis
\item C. Repetundae
\item D. Maiestas
\end{itemize}

\item Given the definition, what is the correct roman term? "The comitia tributa was organized by the location of residence and they vote in an order determined by lot."
\begin{itemize}
\item A. Senatus Consulta
\item B. Plebiscita (singular: plebiscitum )
\item C. Imperium
\item D. Comitia tributa
\end{itemize}

\item Given the definition, what is the correct roman term? "Criminal intent."
\begin{itemize}
\item A. Dolus
\item B. Absolvo
\item C. Condemno
\item D. In iure
\end{itemize}

\item Given the definition, what is the correct roman term? "The power to command held by senior Roman magistrates."
\begin{itemize}
\item A. Imperium
\item B. Legis actiones
\item C. Ius honorarium
\item D. Leges
\end{itemize}

\end{enumerate}

\subsection{Glossary Drill Answer Key}

\subsubsection{Greek Answer Key}

\begin{enumerate}

\item A. Archon Eponymous

\item B. Kyrios

\item A. Boule

\item A. Dike

\item D. Nomothetai

\item C. Strategoi

\item B. Archon

\item A. Paragraphe

\item C. Archon Basileus

\item B. Graphe

\item B. Archon Polemarchus

\item D. Prosklesis

\item B. Logographoi

\item B. The Forty

\item D. Areopagus

\item B. Heliaia

\item D. Apagoge, ephegesis, endeixis, apographe, eisangelia, probole, dokimasia, euthynai

\item C. Eisangelia

\item B. Polis

\item D. The Eleven

\item C. Nomos/Nomoi

\item B. Anakrisis

\item A. Ekklesia

\item C. Miasma

\item A. Apagoge

\item D. Dikastai/dikastes

\item C. Thesmothetai

\item A. Trauma ek pronoias

\item A. Hubris

\item A. Apographe

\end{enumerate}

\subsubsection{Roman Answer Key}

\begin{enumerate}

\item B. Praetor

\item D. Comitia centuriata

\item B. Actus reus

\item B. Ambitus

\item C. Leges

\item B. Stuprum

\item C. Censor

\item A. Mos maiorum

\item C. Intercessio

\item C. Condemno

\item A. Mens rea

\item C. Vis

\item B. Ius honorarium

\item B. Maiestas

\item D. Praetor urbanus

\item A. Ius gentium

\item B. Quaestors

\item D. Aedile

\item B. Calumnia

\item B. Senatus Consulta

\item B. Plebiscita (singular: plebiscitum )

\item B. Lex Aquilia

\item B. Repetundae

\item D. Lex Calpurnia

\item D. Lex Hortensia

\item D. Concilium plebis

\item C. Provocatio

\item A. Paterfamilias

\item A. Lex Iulia de adulteriis coercendis

\item A. Absolvo

\item A. Praetor peregrinus

\item C. Quaestiones perpetuae

\item B. Peculatus

\item A. Delictum

\item C. In iure

\item C. Iudex

\item A. Cognitio extraordinario / Cognitio extra ordinem

\item B. Legis actiones

\item B. Interdictio aquae et ignis

\item D. Comitia tributa

\item A. Dolus

\item A. Imperium

\end{enumerate}
% END GENERATED GLOSSARY DRILL

\end{document}
